\chapter{知识的共享}
\label{chap:ket-sharing}

先看一个很常见的处境：同一张花卉照片，既想用来识别物种，又想用来判断病害，还想标出每朵花的位置。如果为这三件事各训一个独立模型，它们会不约而同地重复学习同一批底层能力——从像素里提取边缘、纹理、形状。既然这些前期计算高度重合，让三个任务共用它们，显然有机会省去大量重复的估计和数据需求，这正是多任务学习最朴素的动机。但“能共用”不等于“能全同”：三件事关注的信息存在微妙的冲突——物种分类认为可以忽略甚至该抹平的细小斑点，恰恰可能是判断病害的关键依据；分类时无关紧要的绝对位置，到了定位任务却必须精确输出。所以本章要解决的，从来不是“怎样把三个模型做成处处相同”，而是一个更精细的问题：怎样充分利用它们的共同规律，同时又为每项预测保留它各自不可或缺的差异。

形式化地说，让多个不同的预测要求共同参与训练，使得一项任务的经验能够改变另一项任务所使用的模型，这种范式称为\term{多任务学习}（multi-task learning, MTL）。它和第13章的迁移学习是一对互补视角：迁移是“老任务帮一个新任务”，多任务则是“多个当前任务一起学、互相帮”。本章以一个常驻系统里的三项任务为例：花卉类别 $\mathcal T_1$、病害判断 $\mathcal T_2$、对象定位 $\mathcal T_3$；推广到一般情形时用 $K$ 表示参与任务数，本例 $K=3$。要说清共享到底有没有好处，前提是先为每个任务约定“错了怎么计量”。设任务 $k$ 的预测器为 $f_k$，损失 $\ell_k$ 给单次预测打分，联合分布 $\mathcal P_k$ 描述它未来面对的输入与标签。在该分布上取平均，就得到我们希望尽量降低的任务风险：
\begin{equation}
  R_k(f_k)=
  \mathbb E_{(\bm{x},y_k)\sim\mathcal P_k}
  [\ell_k(f_k(\bm{x}),y_k)].
  \label{eq:ket-sharing-task-risk}
\end{equation}
这里标签 $y_k$ 在三个任务里分别代表类别、病害状态或位置，所以三者的输出空间和损失单位本就可以各不相同（一个是 0/1 类标、一个可能是患病概率、一个是坐标框）。把 $R_k$ 在训练集 $S_{\mathrm{train},k}$ 上做样本平均，得到经验风险 $\hat R_k$；只对当前一个小批次计算的损失记作 $L_k$。务必区分这三个层次：把训练批次损失 $L_k$ 降下去，只说明模型更贴合这批监督，完全不能证明“共享”降低了未来的真实风险——后者必须用既没参与拟合、也没参与选择的评价样本来估计。本章大量讨论 $L_k$ 怎么合成，但心里要清楚那只是训练端的代理量。

一个完整的共享方案必须同时回答四个问题：谁和谁在一起学、共同使用的到底是什么、各项任务的错误之间怎么取舍、以及愿意为此投入多少标签和算力。本章四个主体节正好对应这四项设计；为了能归因，比较其中某一项时应尽量把其余三项固定。要强调它们可以自由组合、也会彼此牵制，并不是四种只能选一的互斥算法。还可以借此划清本书的分工：第~\ref{chap:ket-transfer-adaptation}章让“已有经验”服务于一个新目标，而本章让“多个当前任务”一起形成经验；倘若各模型都已经分别训练好，你只想把它们的计算或参数组合起来，那属于知识融合，应回到第~\ref{chap:knowledge-fusion}章。

区分共享方式还有一条有用的线索：看联系发生在哪个阶段。三个任务“共用同一批图像”并不等于它们“共用同一套表示”，而共用了表示也不代表它们的预测输出必须满足某种关系——样本、表示、预测这三者，可以分别通过样本组织、前向计算、训练约束这三条不同的路径发生联系。而在训练内部，又要把三个环节分开，因为它们各有各的失效方式：任务权重规定了“想降低哪些错误”，采样和缺失标签处理决定了“实际拿到了什么监督”，梯度协调才把这些监督“翻译成参数的实际变化”。这是一条流水线，后一步无法弥补前一步的缺陷：哪怕梯度协调设计得再精巧，前面若漏掉了某个任务、或标签系统性偏向某类样本，最终更新也不会自动把它补回来。

\section{共享任务的组织}
\label{sec:ket-sharing-organization}

组织共享的第一步不是急着设计网络，而是先厘清“责任”：类别、病害、定位虽然都读照片，但它们在部署时的地位可能完全不同。如果三项输出都要真正交付给用户，那每一项都必须单独验收，谁也不能被牺牲；可如果最终只上线类别识别，那病害和定位完全可以退居幕后，只在训练阶段充当帮助类别学习的“陪练”。先把这种收益要求讲清楚，再去测量任务之间谁帮谁、谁害谁，才能决定到底是三项全在一起训，还是只让其中一部分成组。这里有两个正交的决策：主辅地位回答“要保护谁的收益”，分组与层级回答“谁和谁共享、共享多深”；二者可以同时用来描述一个方案。而且分组不是一锤子买卖，若试训反复表明病害持续受损，就应当依据证据回头修订分组。

\subsection{候选任务与收益要求的确定}
\label{sec:ket-sharing-candidates}

如果终端必须同时交付类别、病害、位置三项，就得为每一项分别设立最低标准——这是硬约束，不是事后平均出来的。要特别避免一个偷懒做法：随手设三个相等的任务权重，训完把结果当成部署要求。病害漏检的代价（漏报疫情）和定位偏差几个像素的代价根本不在一个量纲上，等权重在数学上看起来公平，实际上掩盖了责任差异。除了“绝对不能差于某条线”，还可以提一种相对要求：限制共享训练相对“各任务单独训练”不能退化太多。在相同评价条件和相同预算下，记 $R_k^{\mathrm{single}}$ 为单任务基线风险，$R_k^{\mathrm{joint}}$ 为共同训练得到的风险。沿用式~\eqref{eq:ket-reuse-risk-gain}中正收益表示改善的约定，对必须交付的任务集合$C$施加
\begin{equation}
  \Delta_k=R_k^{\mathrm{single}}-R_k^{\mathrm{joint}}
  \geq-\epsilon_k,
  \qquad k\in C,
  \label{eq:ket-sharing-benefit-constraint}
\end{equation}
其中 $\epsilon_k\geq0$ 是事先约定、愿意容忍的风险增量：共享使风险下降时 $\Delta_k>0$（正收益），上升时 $\Delta_k<0$（受损）；令 $\epsilon_k=0$ 就是最严格的“不许退化”。有两点必须澄清：其一，这条相对约束取代不了绝对最低标准——两个方案可能都烂在合格线之下，但其中一个相对另一个仍有正收益，“比烂赢了”不等于能上线。其二，三个任务损失单位不同，$\Delta_k$ 之间不能直接相加求总收益，跨指标的正规汇总见第~\ref{sec:ket-evaluation-sharing}节。另外，风险是用有限评价样本估出来的，验收时还要考虑估计误差；这个式子写的是一项“要求”，并不是训练算法会自动满足的保证。

反过来，若只部署类别识别，就可以引入一些不打算上线的任务当“陪练”，比如预测花瓣颜色、叶缘形状。这类任务称为\term{辅助任务}（auxiliary task）。它的作用是给训练注入额外监督信号，但评价它好坏的唯一标准是主任务（类别识别）有没有因此受益，它本身不需要成为一个能独立交付的能力。它为什么能帮上忙？物种标签往往很稀缺，但“颜色、叶缘”这类属性标签可能更便宜、更密集；属性监督会迫使共同表示把颜色、叶缘这些判别性特征保留下来，于是类别头就不必只靠有限的物种标签去从零摸索所有区分线索——相当于用更容易获得的监督，替主任务把特征学扎实。Caruana 早在多任务学习的经典工作中，就用“多个任务共用一个隐藏表示、各自接独立的预测读出单元”这一结构刻画了这条影响路径\scite{caruana1997}。

但辅助任务不是越多越好，它也可能帮倒忙：辅助监督逼模型保留的特征，未必是主任务需要的。举个典型例子，跨季节识别物种，理想的表示应当对季节性色差保持容忍（同一朵花在不同光线下仍是它）；可若你加一个“精细预测光照色差”的辅助目标，就会逼着表示去编码、放大这些本该忽略的色差，白白占用共同表示有限的容量。所以一个辅助目标留不留，要看重复训练里主任务收益的均值和波动，再核算标注与部署费用，绝不能用“辅助任务自己的损失降得很好”来代替对主任务的检查。还有一个工程细节：辅助头若只在训练时提供损失，部署阶段可以删掉；但如果主任务的类别预测在计算上真的读取了属性预测的输出（形成了前向依赖），那删掉它就会改变主任务，必须连同这条依赖一起重新验证。

辅助预测问题还可以从已有训练样本构造。Ando和Zhang让多个预测任务共同学习结构，同时允许任务特有函数分别拟合目标；共同结构限制的是各任务可选择的表示或假设空间\scite{ket-ando2005-structure}。构造目标是否有效、标签是否可靠及是否泄漏待预测答案，由第~\ref{part:efficient-knowledge-use}部分讨论。这里先把候选目标与收益标准对应起来，下一步再测量它在共同训练中实际帮助谁。

\subsection{任务间共享关系的估计}
\label{sec:ket-sharing-affinity}

决定谁和谁成组之前，最好先实测任务间的相互影响，而不能只看标签表层的相关性。举例说，物种标签和病害标签确实常常相关，但这背后有两种截然不同的解释：可能是物种本身真的决定了易感病害（这是可利用的规律），也可能只是采样偏差——某些物种恰好只拍到了患病植株（这是虚假相关）。仅凭相关性，你无法判断类别监督会不会反而把病害模型带向一条不可靠的捷径。于是需要一种更直接的“因果式”测量，在已经选定的结构里做个受控实验：从同一个参数状态出发，只单独施加某一项任务的更新，然后观察其他任务的损失到底是升还是降。

把阶段$t$的共享模型参数记为$\bm{\theta}_{\mathrm{s}}^{(t)}$，任务私有模型参数记为$\bm{\theta}_k^{(t)}$。为隔离共享块的作用，固定当前批次和所有私有模型参数，先计算任务$k$的共享梯度$\bm g_k^{(t)}$，再用共同步长$\eta>0$试做一步：
\begin{align}
  \bm g_k^{(t)}
    &=\nabla_{\bm{\theta}_{\mathrm{s}}}L_k^{(t)},
  \\
  \widetilde{\bm{\theta}}_{\mathrm{s},k}^{(t)}
    &=\bm{\theta}_{\mathrm{s}}^{(t)}-\eta\bm g_k^{(t)}.
  \label{eq:ket-sharing-virtual-step}
\end{align}
这一步更新纯粹是“试探”，测完立刻把参数恢复原状，所以称为虚拟更新；测下一个任务时也一律从同一个原始状态出发，保证各次测量互不污染。记 $L_{j\mid k}^{(t)}$ 为“试用了任务 $k$ 的共享块更新之后”，任务 $j$ 在固定批次上的损失。为了让这个变化量和任务 $j$ 自身的量级可比，在 $L_j^{(t)}>0$ 时定义一个有方向的量——\term{任务亲和性}（task affinity）：
\begin{equation}
  A_{k\to j}^{(t)}
  =1-\frac{L_{j\mid k}^{(t)}}{L_j^{(t)}}.
  \label{eq:ket-sharing-affinity}
\end{equation}
这个数的含义很直观：任务 $k$ 的那一步试探更新，让任务 $j$ 的损失相对它原来的水平下降了多大比例——正值代表帮忙（损失下降），负值代表帮倒忙（损失上升）。注意它是有方向的：$A_{k\to j}$ 和 $A_{j\to k}$ 是两回事。Fifty 等人正是把训练全程测得的这种两两影响汇总起来，作为自动搜索任务分组的依据\scite{ket-fifty2021-affinity}。

不过一次虚拟更新的结果很“娇气”，会受当前批次和训练阶段影响。为了让不同方向的亲和性可比，测量时要把步长、共享层范围和评价批次都固定，并在多个批次、多个训练阶段上记录平均值与波动，而不是采信单次读数。还有个数值陷阱：当任务 $j$ 的损失本身已接近零时，分母极小，相对变化会被放大得很夸张，这时应补报绝对变化量，或明确换一种归一化尺度。也要认清批次测量的身份：它给的只是“当前训练动态”下的影响；如果你拿某个验证集去挑选分组，那这个验证集就已经参与了模型选择，最终测试必须另留一份干净数据。

你可能会问：这种影响难道每次都得真的试一步才知道吗？有没有更省事的信号？有——两个任务梯度的内积就能在局部上提示相互影响。道理来自一阶泰勒展开：固定私有参数，把任务 $j$ 的损失沿“任务 $k$ 的更新方向”展开；在共享梯度局部 Lipschitz 连续的条件下，高阶余项是步长的二阶小量，主导项只剩内积：
\begin{equation}
  L_{j\mid k}^{(t)}-L_j^{(t)}
  =-\eta(\bm g_j^{(t)})^{\top}\bm g_k^{(t)}
  +O(\eta^2).
  \label{eq:ket-sharing-local-effect}
\end{equation}
于是在内积非零、步长足够小时，方向由一阶项决定：两个梯度夹角为锐（内积正）时，沿 $k$ 的方向走会顺带降低 $j$ 的损失；夹角为钝（内积负）时则会抬高它。但这里有个关键的不对称来源必须点破：纯内积 $\bm g_j^\top\bm g_k$ 交换 $j,k$ 后不变，是对称的；可一旦步长不是无穷小，结果还要受两个任务各自损失曲率的影响，何况亲和性又除以了各自不同的原损失做归一化。所以“内积对称”完全推不出“实测亲和性也对称”——$k$ 帮 $j$，不等于 $j$ 帮 $k$。

为了亲手验证“内积对称、亲和性却可以不对称”，用一个能完全手算的一维例子。设共享标量 $\theta$ 代表某种特征强度，类别损失 $L_{\text{类}}(\theta)=\tfrac12(\theta-1)^2$ 是个缓坡（曲率小），属性损失 $L_{\text{属}}(\theta)=5(\theta-0.1)^2$ 是个陡坑（曲率大、最优点在 0.1）。在零点处两者梯度恰好都为 $-1$，内积为正；取步长 $\eta=0.3$，无论用哪个任务的更新，提出的新参数都是 $0.3$。走到 $\theta=0.3$ 后，缓坡的类别损失从 $0.5$ 降到 $0.245$（确实受益），但陡坑的属性损失从 $0.05$ 反升到 $0.2$——这一步对它而言迈过了最优点 0.1、冲过头了。于是两个方向算出的亲和性截然不同：$A_{\text{属}\to\text{类}}=0.51$（属性任务的更新帮了类别），$A_{\text{类}\to\text{属}}=-3$（类别任务的更新严重害了属性）。这个非对称的根源是曲率差异和步长过大，而不是内积。作为印证，若把步长缩小到 $0.01$，这一步足够小、谁都不会冲过最优点，两项损失就都会下降。

同一例子还能区分“互相帮助”和“互相干扰”两种几何：当两个正曲率二次损失的最优点都在当前位置的同一侧时，朝那一侧走足够小的一步就能同时降低两者（共赢）；可若把病害损失换成 $\tfrac12(\theta+1)^2$，它在零点的梯度变成 $+1$，与类别梯度 $-1$ 正好要求相反方向的移动——这才是真正的梯度冲突，怎么取正步长都难以同时满足。要记住这些都是在“假设的局部二次损失”上做的计算，是为了建立直觉，并不是在给真实花卉任务的关系下结论；一旦更换表示、标签、共享层，或者只是进入另一个训练阶段，任务关系都可能改变，必须重新测量。

还要把这种测量和第~\ref{sec:ket-reuse-task-relations}节的“顺序迁移”测量区别开。Taskonomy 的做法是先把来源表示固定下来，再在上面训练目标任务的读出头，考察的是“一个现成表示可不可复用”\scite{ket-zamir2018-taskonomy}；而共同训练允许目标任务的梯度反过来改变共享表示本身。这一差别带来一个容易误判的结论：你能从一个类别表示里很好地读出属性（表示里含属性信息），并不代表“把属性训练加进来就一定帮助类别”——信息存在是一回事，联合优化的动态是另一回事。落地时：只部署类别就重点测“属性对类别”这一个方向；三项都交付则要把每一对方向的帮助、受损和估计的不确定性都逐项记录。

\subsection{任务分组与共享层级的确定}
\label{sec:ket-sharing-grouping}

有了亲和性，就可以着手分组。假设实测发现：在当前骨干里，类别和定位的梯度比较相容，而病害持续要求另一组纹理（与前两者频繁冲突）。一个自然的候选是把类别、定位绑成一组共享，病害单独走自己的分支；然后把这个方案和“全部共享”“全部分开”两个极端放在一起比较。但要清醒：亲和性只是用来提出候选分组的线索，不能直接拍板最终方案。原因在于，三个任务一旦真的同训，会同时改变梯度如何合成、容量如何竞争、各项采样机会如何分配——这些“合在一起”的效应，两两测量时根本观察不到，所以两两都互相帮助，也不能保证三个一组整体就好。

比较分组时还必须把模型大小一并算进去，否则结论会被容量差异污染：把一个网络拆成三个同样大的网络，总参数立刻增加约三倍；反过来用一个大骨干替换三个小网络，改变的也不只是共享关系，还白赚了更多容量，性能提升未必来自“共享”。Standley 等人正是在不同网络大小、样本规模和算力预算下系统研究任务分组，说明任务关系的证据必须连同这些条件一起解释\scite{ket-standley2020-grouping}。务实的流程是：先用便宜的局部亲和性测量，从大量可能分组里筛出少数候选；再在相同的总预算下把每个候选配齐、完整训练，按逐任务的验证指标选定方案；最后才在独立测试集上给结论。局部筛选只省下了一部分试训成本，永远不能代替整组的完整验证。

共享的组织单位不一定是“不同的预测要求”，也可以是“同一个要求在不同地区的副本”——比如各个植物园都要做类别识别，但拍摄条件各异。这自然引出层级（hierarchical）共享：用 $a$ 索引植物园，$c(a)$ 表示它所属的地区组。直观想法是把参数拆成三个尺度的贡献之和：所有植物园共享的通用部分（一起学边缘、基础纹理）、同一地区组共享的部分（区域气候、常见设备）、以及单个植物园私有的部分（本地特有的拍摄或校准差异）。需要强调，能写成这种加法分解的前提是参数块同维、且坐标一一对应：
\begin{equation}
  \bm{\theta}_a
  =\bm{\mu}+\bm b_{c(a)}+\bm\delta_a,
  \label{eq:ket-sharing-hierarchy}
\end{equation}
共同中心 $\bm\mu$ 全体通用，组偏移 $\bm b_{c(a)}$ 只在同组植物园间共享，个体项 $\bm\delta_a$ 则为植物园 $a$ 所独有。这里要泼一点冷水：这还只是一个结构性的分解，并没有指定怎么估计；而且若不加额外约束（如正则或先验），同一个总参数可以有很多种等价的拆法，分解并不唯一。地理上的地区分组提供的只是“共享候选”，绝不是正确性保证——两个相邻植物园若用了完全不同的相机系统，按地理硬分组反而可能把观测差异最大的两方绑在一起。

层级共享默认数据能集中；但现实中各植物园可能因数据主权或带宽，无法把训练样本汇集到一处。这时可以让每个节点保留自己的模型，只在训练过程中协调——这就是联邦多任务学习的场景。MOCHA 为每个节点维护模型 $\bm w_a$，同时用一个描述节点间关系的矩阵 $\bm\Omega$ 把它们联系起来，把“本地经验损失”和“跨节点关系正则”放进同一个目标，并交替地更新各节点参数和这个关系结构\scite{ket-smith2017-mocha}。关系矩阵固定后，各节点只用本地数据近似求解自己的子问题，再把必要的统计量上交参与全局协调，最终不必被迫采用同一个预测器——节点之间可以保留差异。不过它那套优化保证是有条件的（针对凸损失、特定正则，且要考虑局部求解误差和节点掉线），不能直接照搬到任意非凸神经网络。还有一个常见误解必须澄清：训练样本不离开节点，并不等于更新里就没有样本信息泄漏——梯度或参数更新仍可能反推出个体数据；真正的隐私保证（如差分隐私）要参考\BookChapterRef{chap:trust-privacy}{卷3《基础、理论与可信性》的“隐私保护”章}。

联邦语境下有三种目标不同的玩法，要分清：一是协调训练出一个大家共用的模型；二是先共同准备一个底座、再各自本地适应；三是联合学习但最终为每个节点保留不同的个性化模型。FedAvg 每轮把各节点本地训练后的参数平均汇总成新的共同模型，属于第一种，其基本流程可参考\BookChapterRef{chap:machine-learning-paradigms}{卷3《基础、理论与可信性》的“机器学习范式”章}；注意“做了本地训练”并不自动等于“最终交付个性化模型”，FedAvg 平均后大家拿到的仍是同一个模型。pFedMe 则明确走第三条路：在共同中心 $\bm\mu$ 附近，为每个节点近似求解一个带近端惩罚的本地目标 $L_a(\bm w_a)+\lambda\|\bm w_a-\bm\mu\|_2^2/2$，再通过协调去压低这些“局部最优解”的平均。式中 $L_a$ 是节点 $a$ 的本地平均损失，$\lambda>0$ 像一根橡皮筋，把个性化模型拴在共同中心附近、控制其偏离幅度，最终各节点可以保留各自不同的 $\bm w_a$。这种“近端联系”给出了另一种共享的目标函数，但它同样需要分别核查本地求解误差、节点如何采样、通信开销以及逐节点的真实收益；尤其不能用“共同模型的平均成绩不错”去冒充“每个节点的个性化效果都有保证”——平均会掩盖个别节点的严重退化\scite{ket-dinh2020-pfedme}

最后回到“任务到底以什么为单位”，这在指令式、异构数据训练里尤其容易含糊。同一个预测要求（比如“判断是否患病”）被改写成十个表述模板，并不自动变成十项独立能力——它们本质是同一项任务的不同说法；反过来说，同一个模板翻译成另一种语言，也不能假装语言和数据来源的差异不存在。正确做法是分别记录“语义要求、数据来源、语言、模板”这几个维度，才能合理决定怎么分组、评价时按什么单位汇总。等确定了哪些单位之间可以交换经验，下一节再回答交换“发生在哪里”——是共同计算里、参数关系里，还是预测输出关系里。

\section{可共享知识的设计}
\label{sec:ket-sharing-structure}

定好谁和谁一组之后，下一个设计问题是：这些任务到底在网络的哪个位置共享？要先破除一个误解——同组任务并不需要处处使用同一个网络。共享可以发生在深浅很不相同的三个位置：类别和病害可以共用一段特征提取（共同计算）；也可以前向计算完全独立、只让各自的参数估计受到同一个统计先验约束（参数关系）；定位的输出还可以直接当作病害预测的输入、或仅在训练时对病害出现的位置施加约束（输出关系）。本节就按这三个位置依次展开。还要提醒，若一个“共享表示”本身就是由“共享参数”算出来的，那它们只是同一条共享路径的两种说法，不要当成两类互不相干的方法重复计数。

在让样本、表示或预测被多个任务共用之前，有一个看似琐碎却致命的前提：必须确认接口两端指向的是同一个对象。如果图像裁剪方式不同、坐标原点不一致、或者标签和图像的对应关系错位，那么即便两个张量形状一模一样，它们也无法交换——形状相同不等于语义相同。有两类经典的“统一接口”工程值得借鉴，也需要警惕：T5 用任务前缀加文本答案，把五花八门的要求都接进同一个“文本到文本”接口\scite{raffel2020t5}；MultiModel 则为图像、语音、文本各自保留观测编码器和预测读出模块，只让它们共用内部的核心计算\scite{ket-kaiser2017-multimodel}。但要清醒，统一的只是形式接口：它既没有统一各任务的难度和数据规模，也不代表原论文的所有实验都用了“同时联合训练”。判断这类接口好不好，仍要看它是否为当前每个任务保留了所需的特征。

\subsection{共同计算结构的设计}
\label{sec:ket-sharing-computation}

如果多个任务共用一段映射，那每个任务的损失都会反过来改变这段映射，共享设计于是变成三个具体选择：共享到哪一层为止、在哪些分支之间交换特征、以及这种连接或混合比例是否随每次输入而变化。这里有一条重要的分界线：固定连接可以在训练中学会一组固定系数，但部署后对所有输入都一样；条件连接则在使用时根据任务身份甚至当前输入内容动态选择——不过只依赖任务身份的选择，对同一个任务而言仍是固定的。下面用“固定连接”和“逐输入混合”作为两大对照来展开，两者都允许任务保留自己的私有分支。

\subsubsection{固定共享范围的设计}
\label{sec:ket-sharing-fixed}

最简单的固定共享结构是硬共享：对同一张照片，只算一次共同特征，再由几个不同的“读出部件”分别给出类别、病害和位置。通常把前面那段共同映射称为骨干（backbone），把后面对应每项预测的映射称为任务头（task head），计算流程为
\begin{equation}
  \bm z=h_{\bm{\theta}_{\mathrm{s}}}(\bm x),
  \qquad
  \hat y_k=q_{\bm{\theta}_k}(\bm z).
  \label{eq:ket-sharing-backbone}
\end{equation}
这里 $\bm z$ 是所有任务共用的表示，由共享参数控制的骨干 $h_{\bm{\theta}_{\mathrm{s}}}$ 算出；任务头 $q_{\bm{\theta}_k}$ 用各自的私有参数把 $\bm z$ 读成预测 $\hat y_k$。训练时，三个任务的损失都沿自己的头回传、汇聚到同一个骨干，相当于逼着骨干同时满足多项要求。共享多深是个需要权衡的决策：只共享浅层时，各任务的深层分支可以分别保留自己需要的空间细节；共享更多层能省下重复计算，却要求同一个高层表示同时装下所有任务必需的特征，而这些特征可能互相打架。一个典型矛盾是：分类乐于使用对位置不敏感的汇总特征（花朵在哪无所谓），定位却必须保住特征图上的空间对应关系。正因为这种差异，骨干该共享多深没有通用答案，只能靠具体任务的实验来定。

“共享表示”并不意味着各任务必须用同一个预测函数——它们可以共享一组“积木”，但用不同方式搭。为看清这一点，先用同一 $d$ 维输入上的标量线性预测做简化：让 $\bm U\in\mathbb R^{d\times r}$ 提供 $r$ 个共同的方向（一组共享基），任务 $k$ 再用自己的系数 $\bm a_k\in\mathbb R^r$ 去挑选和组合这些方向。将各任务模型参数并列成矩阵，便有
\begin{equation}
  \hat y_k=\bm a_k^{\top}\bm U^{\top}\bm x,
  \qquad
  \bm W=\bm U\bm A,
  \label{eq:ket-sharing-subspace}
\end{equation}
把各任务的系数并成矩阵 $\bm A$（第 $k$ 列是 $\bm a_k$），则所有任务的参数可紧凑写成 $\bm W=\bm U\bm A$（$\bm W$ 的第 $k$ 列即任务 $k$ 的参数向量）。这样每个任务的参数方向可以不同，但都被约束在 $\bm U$ 张成的那个低维子空间里——共享的是“允许活动的方向集合”。经典的多任务特征学习在此基础上再加一层“系数行稀疏”约束，鼓励所有任务只使用少数几个共同特征（让某些共同方向整体被关掉）；它和这里预先固定秩 $r$ 的做法有关联、都在压缩自由度，但并不是同一个约束\scite{ket-argyriou2006-features}。训练联合估计共同结构与任务系数，部署按选定映射计算预测。限制太强会排除任务所需方向，限制太弱则可能不再有效减少估计自由度。

硬共享（完全共用骨干）若限制太强，还有一个折中：保留各自的分支，只在指定的某几层让分支之间交换一点信息。Cross-stitch（十字绣）网络用一个可训练的小矩阵，去线性混合那些形状已经对齐的分支特征。以两个分支在同一空间位置、同一通道上的标量激活 $h_1,h_2$ 为例，混合后供两个分支继续使用的值为：
\begin{equation}
  \begin{bmatrix}
    \widetilde h_1\\ \widetilde h_2
  \end{bmatrix}
  =
  \begin{bmatrix}
    a_{1,1}&a_{1,2}\\a_{2,1}&a_{2,2}
  \end{bmatrix}
  \begin{bmatrix}
    h_1\\h_2
  \end{bmatrix}.
  \label{eq:ket-sharing-cross-stitch}
\end{equation}
这个矩阵的对角项负责缩放本分支自己的激活，非对角项则负责把对方分支的信息“掺”进来——掺多少由训练自动学。要注意这些系数是可正可负的实数，不是概率，行和也不要求为 1，所以别把它当注意力权重解读。上式只展开了两个分支在单个位置上的混合；任务更多时，可以在事先选定的若干分支之间设置对应的混合矩阵\scite{ket-misra2016-cross-stitch}。

举个数字例子：若学出的混合矩阵两行为 $(0.9,0.1)$ 和 $(0.2,0.8)$，某位置两支激活为 $2,6$，则混合后第一个分支拿到 $0.9\times2+0.1\times6=2.4$，第二个拿到 $0.2\times2+0.8\times6=5.2$。反向传播会同时更新混合系数和各分支原参数。它的“固定”性质要强调：部署后这组系数就是训练好的常数，不会因为换了张照片而由某个网络重新计算；照片改变的只是被混合的激活数值本身，混合比例不变。连接在哪一层、哪些分支参与，也都是事先固定的——这正是它和后面“按输入动态门控”的本质区别。

至此骨干—头、共享子空间、Cross-stitch 三种固定结构，本质都在同一条权衡曲线上取舍：共享得少，任务自由，但可能重复学习同样的纹理、浪费数据和算力；共享得多，省估计，却可能挤不出足够的方向去同时保留病害细节和位置信息。因此在比较“共享几层”“在哪交换”时，一定要控制总参数量、特征维数和算力预算相同——否则一个更宽的分支单凭容量更大带来的收益，会被你错误地记到“共享机制更好”的账上。

\subsubsection{条件共享范围的设计}
\label{sec:ket-sharing-conditional}

固定的共享范围即使能区分任务，也未必能应对同一任务内部的变化。同样是病害判断：叶片清晰时，模型应主要依赖斑点纹理；可当叶片被遮挡时，又得更多依靠整体形状和上下文。这种“同任务、不同输入、需要不同专家”的情形，靠一组固定混合比例是处理不了的，但为每种情形各训一个完整模型又太浪费。自然的折中是：保留多种可复用的映射，让当前输入自己决定此刻该以什么比例组合它们。

\term{多门控混合专家模型}（Multi-gate Mixture-of-Experts, MMoE）把这些可复用映射称为“专家”，并为每一项任务单独配一个门控网络，由它根据当前输入算出该用各专家的比例。设有 $E$ 个专家 $\bm e_m(\bm x)$，每个都输出同维的表示向量；任务 $k$ 的门控打出原始分数 $a_{k,m}(\bm x)$，经 softmax 变成非负、和为 1 的混合权重，再对专家向量加权求和：
\begin{align}
  g_{k,m}(\bm x)
    &=\frac{\exp(a_{k,m}(\bm x))}
      {\sum_{r=1}^{E}\exp(a_{k,r}(\bm x))},
  \\
  \bm z_k(\bm x)
    &=\sum_{m=1}^{E}g_{k,m}(\bm x)\bm e_m(\bm x),
  \nonumber\\
  \hat y_k&=q_{\bm{\theta}_k}(\bm z_k(\bm x)).
  \label{eq:ket-sharing-mmoe}
\end{align}
其中 $g_{k,m}$ 就是在当前输入上专家 $m$ 对任务 $k$ 的权重，$\bm z_k$ 是该任务混合得到的专属表示，最后由自己的头 $q_{\bm{\theta}_k}$ 读成预测。MMoE 同时提供了两种灵活性：不同任务对同一组专家可以用不同比例（任务间差异），同一任务在不同输入下比例也能变（任务内、随输入的差异），从而比固定共享更细致地表达“何时该共享什么”\scite{ket-ma2018-mmoe}。

训练动态上，任务 $k$ 的损失会沿自己那条路径更新它的头和门控，同时也更新被它用到的专家，于是专家逐渐学成“被若干任务以不同方式调用”的通用部件；预测时门控和专家表示都随当前输入重新计算。可以用两个退化对照看清它的定位：若让所有任务共用同一个只看输入的门控，那么同一张图对所有任务给出的混合表示都相同（丢掉了任务间差异）；若换成 Cross-stitch，则混合系数根本不随输入重算（丢掉了任务内差异）。三者能复用相似部件，但“凭什么来选择”这件事各不相同。

需要保留地区或任务差异时，还可在专家以外增加私有分支或校准模型参数。新增自由度也需要训练样本来估计，结构本身不保证个性化收益。样本多、更新频繁的任务可能先塑造共同专家，其他任务在有限预算内未必能充分调整；专家数量增加也可能只得到多份近似映射。

MMoE 还容易给人“专家自动学会分工”的错觉，这需要实证检验而不是想当然：应当按任务、按输入子群去统计各门控权重的使用比例和集中程度（是不是所有门控都只依赖同一个专家？），再做删去或替换某个专家的干预实验，看哪些预测随之改变，并与同等预算的固定共享结构公平比较。仅仅看到门控权重不同，不能就给某个专家贴上“颜色专家”之类的语义标签。另外注意，标准 MMoE 是稠密软混合，一次请求通常仍要把所有 $E$ 个专家都算一遍，省的是存储意义上的共享、未必省单次计算；要真正减少单次计算量，得引入稀疏选择和能跳过未选中专家的执行机制，相关预算核算见第~\ref{sec:ket-sharing-budget}节。本节专家由任务共同训练，第~\ref{sec:ket-integration-conditional-modules}节则组合已经形成的模块。

\subsection{共享先验与参数关系的设计}
\label{sec:ket-sharing-prior}

前面的共享都发生在前向计算里。还有一种更“隐式”的共享：哪怕病害模型完全不使用类别模型的前向通路，在它样本稀少、参数估不准时，仍可以借助“别的任务的参数大概长什么样”这一统计信息，来约束自己的自由度、避免在少数据上乱拟合。这时被共享的不是计算，而是参数可能取值的统计关系。讨论它之前必须先固定比较坐标：下面涉及的任务参数块 $\bm w_k\in\mathbb R^d$ 必须同维且坐标一一对应（比如同一组特征上的线性权重，或已对齐、按相同顺序展开的编码器块）。类别、病害、位置那些结构异构的头仍各自保留，绝不能把它们直接拉来求平均；为简洁，后面的公式都省略这些私有部分。

最简单的参数关系假设是：各任务的参数块都围绕同一个共同中心上下波动。数据多的任务主要负责把这个中心估准，数据稀少的任务则在“贴合自己那点观测”和“不要离群体中心太远”之间做加权取舍——中心相当于免费提供了一个强正则。

为把“围绕中心”写成概率模型，可采用下列简化层次高斯先验：
\begin{equation}
  \bm\mu\sim\mathcal N(\bm0,\kappa^2\bm I),
  \qquad
  \bm w_k\mid\bm\mu\sim
  \mathcal N(\bm\mu,\tau^2\bm I).
  \label{eq:ket-sharing-hierarchical-prior}
\end{equation}
式中 $\bm I$ 是 $d$ 维单位阵，两个固定尺度各管一件事：$\kappa$ 限制中心 $\bm\mu$ 本身离零能有多远（给中心也套个正则），$\tau$ 限制各任务参数离中心能散多开（$\tau$ 越小任务被捆得越紧）。模型假设给定中心后各任务参数条件独立、给定参数后各任务样本条件独立，于是联合后验可以干净地因式分解。以 $\mathcal L_k(\bm w_k)$ 记任务 $k$ 的负对数似然，最大化后验就等价于把“数据的负对数似然”和“先验的负对数密度”加在一起最小化：
\begin{equation}
  \sum_{k=1}^{K}\mathcal L_k(\bm w_k)
  +\frac{1}{2\tau^2}\sum_{k=1}^{K}
     \|\bm w_k-\bm\mu\|_2^2
  +\frac{1}{2\kappa^2}\|\bm\mu\|_2^2.
  \label{eq:ket-sharing-hierarchical-map}
\end{equation}
忽略只依赖固定尺度的常数后，高斯先验恰好留下两个二次惩罚：任务参数偏离共同中心要受罚（鼓励聚拢），中心离零太远也要受罚（防止中心整体漂走）。一个很有启发的结果是：固定各任务参数 $\bm w_k$，对中心 $\bm\mu$ 求最优（令导数为零），得到
\begin{equation}
  \bm\mu
  =\frac{\sum_{k=1}^{K}\bm w_k}
    {K+\tau^2/\kappa^2}.
\end{equation}
注意分母里的 $\tau^2/\kappa^2$：正因中心自己也被 $\kappa$ 往零方向拉，所以最优中心并不是各任务参数的简单算术平均（只有 $\kappa\to\infty$、中心不被约束时才退化成平均）。中心定下来后，每个任务仍要努力拟合自己的数据，只是数据少、似然约束弱的任务会相对更听先验的话、被拉得更靠近中心——但它是被“拉”，不是被直接“设成”中心，自身观测仍在起作用。当然，这套机制能否带来好处，根本上取决于“存在单一共同中心”这个假设是否合理。

如果任务天然分成几群（比如“识别类”任务和“定位类”任务参数规律不同），单一中心就太粗糙了。可以引入簇身份 $c_k$，让任务 $k$ 的参数围绕所属簇的中心 $\bm\mu_{c_k}$ 变化，并顺带估计每个任务属于各簇的概率。Bakker 和 Heskes 用贝叶斯多任务模型系统研究了共同先验、任务聚类和门控三者的联系\scite{ket-bakker2003-clustering}；上面的单中心模型只演示了“向中心收缩”这一核心机制，并不复现该论文完整的推断。簇的来源有两种：可以依据已经核验过的任务关系直接指定，也可以交给数据去估；后者需要足够样本才能把不同簇区分开，不能只因为“加上聚类后拟合更灵活、训练损失更低”就断言分簇是真实可靠的。

共同中心的假设只管“大家都别离一个点太远”，它刻画不了更细的关系：哪些任务的参数倾向一起变大、哪些甚至倾向朝相反方向变化？\term{任务关系学习}（Multi-Task Relationship Learning, MTRL）改用任务参数之间的协方差矩阵来表达这种“联动”关系。把各任务参数并成矩阵 $\bm W=[\bm w_1,\ldots,\bm w_K]\in\mathbb R^{d\times K}$，它的每一行对应同一个参数坐标在 $K$ 个任务上的取值，于是“任务两两如何协变”的信息正好可以放在一个 $K\times K$ 的协方差里。具体采用一个矩阵正态先验：$\bm W$ 的各行相互独立，每一行都是零均值、协方差同为 $\bm\Omega\succ0$（正定）的高斯向量，即同一参数坐标在各任务上的取值共享这个任务间协方差。写出它的负对数密度、丢掉与未知量无关的常数，剩下的关键部分为：
\begin{equation}
  \frac12\operatorname{tr}
    (\bm W\bm\Omega^{-1}\bm W^{\top})
  +\frac d2\log|\bm\Omega|.
  \label{eq:ket-sharing-matrix-prior}
\end{equation}
这里 $\operatorname{tr}$ 是矩阵的迹，$|\bm\Omega|$ 是行列式。两项各有角色：第一项 $\operatorname{tr}(\bm W\bm\Omega^{-1}\bm W^\top)$ 按任务协方差去度量参数变化——两个强正相关的任务若参数一起变动，代价就小；第二项含 $\log|\bm\Omega|$，来自高斯分布的归一化因子。一个容易犯的错误是：当 $\bm\Omega$ 本身也要从数据估计时，第二项绝不能当常数扔掉，因为它正好起到“阻止协方差无限膨胀”的作用（否则只需把 $\bm\Omega$ 放大就能让第一项趋于零）。为了把关系矩阵的尺度钉死又不引入口味，MTRL 改用“迹等于 1”的约束，构造如下正则化准则：
\begin{align}
  \min_{\bm W,\bm\Omega}\quad&
  \sum_{k=1}^{K}\hat R_k(\bm w_k)
  +\frac{\lambda_1}{2}\|\bm W\|_{\mathrm F}^2
  \nonumber\\
  &+\frac{\lambda_2}{2}
    \operatorname{tr}(\bm W\bm\Omega^{-1}\bm W^{\top})
  \nonumber\\
  \text{满足}\quad&
  \bm\Omega\succeq0,
  \qquad \operatorname{tr}(\bm\Omega)=1.
  \label{eq:ket-sharing-mtrl}
\end{align}
式中 $\lambda_1,\lambda_2>0$ 分别控制参数整体大小和关系惩罚强度，$\|\cdot\|_{\mathrm F}$ 是 Frobenius 范数。迹为 1 的约束替代了原来的对数行列式项，防止靠放大协方差作弊；半正定约束 $\bm\Omega\succeq0$ 则保证它在数学上能当协方差用。这个问题有个很好的性质：关系正则部分对 $\bm W,\bm\Omega$ 是联合凸的，所以当每个任务的经验风险对自己的 $\bm w_k$ 也凸时（典型如固定特征上的线性模型加平方损失），整个约束问题就是凸的、可全局求解。但若把同一套正则搬到非凸神经网络上，就不能只凭正则项凸就声称整体目标凸。还要注意这是一个“改造过的估计准则”，并不是简单删掉行列式项后还能和高斯 MAP 画等号\scite{ket-zhang2010-mtrl}。

这个联合问题通常交替求解：关系矩阵 $\bm\Omega$ 固定时，它退化成各任务带耦合正则的凸子问题，可解 $\bm W$；参数 $\bm W$ 固定后，再反过来更新 $\bm\Omega$。后者有一个漂亮的闭式解，值得推一遍，因为它说明“任务关系”确实是从参数里学出来的。令 $\bm B=\bm W^{\top}\bm W$（一个记录任务两两参数内积的矩阵），$\bm B^{1/2}$ 表示把它的非负特征值逐个开平方所得的半正定矩阵。要严格证明下面的更新在所有可行关系矩阵里最优，思路是先在正定内部建立一个下界，而且不预先假设 $\bm\Omega$ 和 $\bm B$ 特征方向相同。以同形实矩阵的$\langle\bm A,\bm C\rangle_{\mathrm F}=\operatorname{tr}(\bm A^{\top}\bm C)$为Frobenius内积，Cauchy--Schwarz不等式及$\operatorname{tr}(\bm\Omega)=1$给出
\[
  \begin{aligned}
    (\operatorname{tr}\bm B^{1/2})^2
      &=\left\langle\bm\Omega^{-1/2}\bm B^{1/2},
                      \bm\Omega^{1/2}\right\rangle_{\mathrm F}^{2}\\
      &\leq\|\bm\Omega^{-1/2}\bm B^{1/2}\|_{\mathrm F}^{2}
             \|\bm\Omega^{1/2}\|_{\mathrm F}^{2}\\
      &=\operatorname{tr}(\bm B\bm\Omega^{-1}).
  \end{aligned}
\]
最后一行 $\operatorname{tr}(\bm B\bm\Omega^{-1})$ 正是优化准则里要被最小化的那个关系项，所以我们得到了它的一个下界 $(\operatorname{tr}\bm B^{1/2})^2$。当 $\bm W\ne\bm0$ 时，Cauchy–Schwarz 取等号要求存在正常数 $c$ 使 $\bm B^{1/2}=c\bm\Omega$（两者方向重合、只差尺度）；再用 $\operatorname{tr}\bm\Omega=1$ 把尺度定死，便得到最优关系矩阵：
\begin{equation}
  \bm\Omega^\star
  =\frac{\bm B^{1/2}}
    {\operatorname{tr}(\bm B^{1/2})}.
  \label{eq:ket-sharing-mtrl-omega}
\end{equation}
直观地读这个解：若 $\bm B$ 正定，该 $\bm\Omega^\star$ 直接达到下界、即为最优；它和 $\bm B^{1/2}$ 同向，只是被归一化到迹为 1。若 $\bm B$ 有零特征值，最优解落在半正定的边界上（允许某些任务组合“无关”），在满足后面值域条件下，于非零特征子空间内套用同一不等式，下界依然成立。把解代回可知关系项的值就是 $(\operatorname{tr}\bm B^{1/2})^2$，确达到下界。还有个很形象的解释：若 $\bm B$ 的特征值为 $b_j$，那么第 $j$ 个关系方向分到的“权重份额”是 $\sqrt{b_j}/\sum_r\sqrt{b_r}$——沿参数变化越大的方向，在关系矩阵里占比越高，而且是按平方根分配。如此交替更新参数与关系，相当于让当前这批任务共同投票决定正则化所用的关系结构。部署时各任务仍只用自己的参数预测，不必同时加载其他任务模型。

这里补一个严谨性细节：准则允许 $\bm\Omega$ 半正定（可有零特征值，以表达“某些任务方向完全无关”），但式子里写的是普通逆 $\bm\Omega^{-1}$，到了奇异边界上就得小心。正定内部直接用普通逆；落到边界时要改用伪逆，并附加一个值域条件——要求 $\bm W^{\top}$ 的列空间包含在 $\bm\Omega$ 的非零值域内，否则该正则项定义为无穷大（相当于禁止这种搭配）。数值实现时若加一个很小的正定下界把它强行变正定，其实改变了可行集合，闭式更新也要相应调整。还有个退化情形：当 $\bm W=\bm0$ 时，任何迹为 1 的关系矩阵都给出零惩罚，此时绝不能再套上面的闭式解去做除以零的运算。

协方差框架的一个强大之处是它允许任务之间“负相关”——但务必把“统计上的负相关”和“训练上的互相损害”区分开，这是两回事。为看清这点，对同一特征坐标上的两个任务，用一个参数 $\rho$ 控制参数相关性，取
\begin{equation}
  \bm\Omega=\frac12
  \begin{bmatrix}1&\rho\\\rho&1\end{bmatrix},
  \qquad |\rho|<1.
\end{equation}
条件$|\rho|<1$使该矩阵正定且迹为1。将其逆代入关系惩罚，可以直接比较同向与反向模型参数的代价：
\begin{equation}
  \operatorname{tr}(\bm W\bm\Omega^{-1}\bm W^{\top})
  =
  \frac{2\bigl(\|\bm w_1\|_2^2+\|\bm w_2\|_2^2
  -2\rho\bm w_1^{\top}\bm w_2\bigr)}
  {1-\rho^2}.
  \label{eq:ket-sharing-negative-correlation}
\end{equation}
把逆矩阵代进惩罚项就能直接比较：当 $\bm w_2$ 接近 $-\bm w_1$（两个任务参数几乎反向）时，在同样的 $|\rho|$ 下，取负的 $\rho$ 比取正的 $\rho$ 给这组参数更小的惩罚。也就是说，协方差模型能识别并鼓励“互补”关系：如果两个任务描述的是同一对象的互补状态（一个往正方向、一个往负方向），知道其中一个任务的参数，反而有助于约束另一个任务的反向趋势，这是有益的。再强调一次术语：这里的“负相关”描述的是一种先验统计联系；而“负迁移”指的是共同训练相比单任务基线抬高了风险。两者只是名字里都有“负”，含义完全不同，不能划等号。

参数协方差方法还有个根本性的局限：它依赖表示坐标。对某个任务，如果你把一个特征的符号翻转并同步翻转对应参数，或者重排隐藏单元顺序但保持网络预测函数不变，参数向量的表观相关性就可能随之改变——可模型本质没变。所以在参数未对齐时，不能仅凭协方差数值小就断言两个任务无关。协方差刻画的是当前坐标下的统计联系，既不是永久不变的语义关系，更不是因果关系；错误的先验形状、错误的簇结构、或样本不足，都可能导致有害的“不当收缩”（把本该不同的任务硬拉近）。最后划清与第~\ref{sec:ket-transfer-prior}节的分工：那里是先用历史任务备好先验、再交给新任务（迁移），本节是多个当前任务一边共同估计关系、一边立刻使用它（共享），收益必须在这些当前任务上分别验证。

\subsection{共享输出关系与约束的设计}
\label{sec:ket-sharing-outputs}

除了共享计算和共享参数关系，还有第三种更直观的共享位置：直接在任务的预测输出之间建立联系。比如病害分支该重点看哪个区域，定位分支预测出的花朵位置就能直接给它线索——这种联系比绕到参数相关性上去推断要直接得多。但输出层面的共享有严格的对齐前提：两个输出必须指向同一个对象、同一套裁剪坐标。一丛多朵花的照片里，左边那朵的类别不能用来约束右边那朵的病害；叶片压根不在画面里时，没预测出叶斑也不能反推“健康”。只有这些对象和坐标条件先讲清楚，才谈得上两种用法：要么让一个任务的输出真正进入另一个任务的前向计算，要么只让各输出在训练时共同满足某种一致性约束。

第一种用法是预测传递（prediction passing）：让中间任务的预测成为后续任务计算图的一部分。PAD-Net 是代表性工作，它先对图像预测一批中间任务——深度、表面法向、轮廓、语义分割等，再把这些预测编码、汇合成特征，回过头来支持最终的深度和语义预测；无论训练还是测试，这些中间量都由当前图像实时生成\scite{ket-xu2018-padnet}。

为说明花卉任务怎样传递预测表示，另设一个教学计算图，而不把它当作PAD-Net实验或默认接在MMoE之后。共同编码器$h$产生$\bm z$，$t_k$为接收任务形成基础表示，$u_j$给出中间任务预测，随后将其他任务的预测表示加入当前表示：
\begin{align}
  \bm z&=h(\bm x),
  \qquad \bm z_k=t_k(\bm z),
  \nonumber\\
  \widetilde y_j&=u_j(\bm z),
  \qquad
  \bm r_j=\phi_j(\widetilde y_j),
  \nonumber\\
  \bm v_k&=\bm z_k+
     \sum_{j\ne k}
     \bm a_{k,j}(\bm r_k,\bm r_j)
       \odot\psi_{k,j}(\bm r_j),
  \nonumber\\
  \hat y_k&=q_k(\bm v_k).
  \label{eq:ket-sharing-output-message}
\end{align}
其中 $\phi_j$ 把中间任务的预测 $\widetilde y_j$ 编码成“消息”特征 $\bm r_j$，$\psi_{k,j}$ 负责把发送方的特征对齐到接收方的坐标和尺度，门控 $\bm a_{k,j}$ 再逐位置决定接收多少。这里的维度约束很实在：设接收网格高、宽、通道为 $H_k,W_k,C_k$，那么 $\bm z_k$、$\bm v_k$ 以及参与逐元素相乘的两个消息张量，都必须落在同一个 $\mathbb R^{H_k\times W_k\times C_k}$ 空间里。一个图像级的类别向量要参与空间融合，得先映射成通道描述再广播到每个位置；一张定位热图则必须按同一套裁剪坐标重采样到接收网格。再次强调：仅仅把张量 pad 成相同形状，并不能替代真正的坐标对齐。

$\bm a_{k,j}$调节每个位置接收的消息，$\odot$为逐元素乘法，$q_k$把加入消息的$\bm v_k$转成最终预测输出。这样病害分支可利用定位热图与类别线索。训练以中间和最终监督共同约束这些可微路径；部署由预测输入图像重算中间任务预测，再完成依赖它的任务。

预测传递有个必须正视的代价：部署时下游任务拿到的永远是上游的“预测值”，而不是训练时才有的真实辅助标签——不能把真实标签当成一种部署期可免费获得的观测。训练时用真实区域、部署时却只能用预测区域，下游实际面对的误差分布和训练是不一样的（这叫训练—测试输入差异）。用概率图而非硬标签能多保留一些类别概率的分配信息，但这些概率本身准不准并无保证；因此应当专门构造测试，考察“区域偏移、上游给出错误类别、低置信输入”分别会怎样拖累下游。同理，消息门控或注意力权重在未经校准之前，也不能想当然地解释成“这个位置有多可信”的概率。

第二种用法不引入任何部署期依赖：如果你不希望线上预测时各个任务互相等待、互相拖累，可以只在训练损失里加一项，惩罚它们彼此不一致，训练完各任务仍独立运行。为便于说明空间关系，先把病害任务从“图像级有没有病”扩展成“逐像素的病变预测”，并让它和同一对象的区域标注对齐。病变概率$p_{\text{病}}(u\mid\bm x)$和区域概率$p_{\text{区}}(u\mid\bm x)$都定义在非空像素集合$\mathcal U$上；标签若规定病变只能位于该区域内，可以用区域外病变的平均软惩罚：
\begin{equation}
  L_{\mathrm{rel}}
  =\frac{1}{|\mathcal U|}
    \sum_{u\in\mathcal U}
    p_{\text{病}}(u\mid\bm x)
    \bigl[1-p_{\text{区}}(u\mid\bm x)\bigr],
  \label{eq:ket-sharing-output-constraint}
\end{equation}
被求和的乘积在“病变概率高、区域概率低”的位置取大值，正好抓住了“病变跑到了区域外面”这种不一致；除以像素总数取平均，是为了不让物体大小（像素多寡）直接决定惩罚的尺度。在各任务原有监督损失上再加 $\lambda_{\mathrm{rel}}L_{\mathrm{rel}}$（$\lambda_{\mathrm{rel}}\geq0$）就能抑制这种不一致。它有几个明确的适用边界：只输出图像级病害状态的模型用不了它（没有像素对应）；这个约束本身不能替代标签监督，因为“病变预测全为零”或“区域预测全为一”都能轻易让它为零——模型可以靠塌缩来作弊；此外，若区域本身定义错了、或惩罚权重过大，反而会压制掉正确的病变预测。

两种输出级共享的取舍到此清楚：预测传递在部署期增加了前向依赖，也因此会把上游的错误传播给下游；一致性正则只在训练阶段改变优化目标，部署时各任务可以保持独立计算，但代价是它并不保证预测在推理时严格满足那条规则（只是被鼓励在训练数据上一致）。注意与第~\ref{sec:ket-integration-fixed-modules}节的分工：那里讲的是把已经学好的模块串起来，这里讲的是关系如何在共同训练中持续施加约束；至于这些规则从哪获取、怎么核验，则属于监督信号获取的话题。读图~\ref{fig:ket-sharing-locations}时可用一个窍门：实线代表部署时仍要真实执行的前向计算，虚线代表只在训练时存在的先验联系或统计约束——据此就能判断哪些依赖在上线后依然跑、哪些只在训练里起作用。

\begin{figure*}[htbp]
  \centering
  % Three sharing locations and a conditional-computation inset.
% Conceptual diagram only; all dimensions are in millimetres.
\begin{tikzpicture}[
  x=1mm,y=1mm,
  font=\figurefont\footnotesize,
  text=BookInk,
  >={Latex[length=1.8mm,width=1.3mm]},
  panel/.style={draw=BookRule,line width=0.6pt},
  model/.style={draw=BookTeal,fill=FigureModel!10,line width=0.8pt,
    minimum height=10mm,text width=11mm,align=center,inner sep=1mm},
  prior/.style={draw=BookBlue,fill=FigureInput!10,line width=0.8pt,
    minimum height=11mm,text width=27mm,align=center,inner sep=1mm},
  output/.style={text width=14mm,align=center,inner sep=0pt},
  relation/.style={draw=BookAmber,fill=FigureLoss!10,line width=0.8pt,
    minimum height=10mm,text width=40mm,align=center,inner sep=1mm},
  flow/.style={draw=BookInk,line width=0.8pt,-{Latex}},
  wire/.style={draw=BookInk,line width=0.8pt},
  stat/.style={draw=BookBlue,dashed,line width=0.8pt},
  constraint/.style={draw=BookAmber,dashed,line width=0.8pt},
  heading/.style={anchor=west,font=\figurefont\small}
]
  \path[use as bounding box] (0,5) rectangle (169,-160);
  \node[heading] at (1,0) {(a) 表示/模型参数 × 前向计算};
  \node[heading] at (59,0) {(b) 模型参数 × 训练先验};
  \node[heading] at (117,0) {(c) 预测输出 × 训练约束};
  \draw[panel] (0,-7) rectangle (53,-74);
  \draw[panel] (58,-7) rectangle (111,-74);
  \draw[panel] (116,-7) rectangle (169,-74);

  \node[model,text width=27mm] (shared) at (26.5,-20)
    {共同映射\\$h_{\bm{\theta}_{\mathrm s}}$};
  \node[model] (head1) at (8.5,-43) {$q_1$};
  \node[model] (head2) at (26.5,-43) {$q_2$};
  \node[model] (head3) at (44.5,-43) {$q_3$};
  \node[output] (outa1) at (8.5,-62) {类别\\$\hat y_1$};
  \node[output] (outa2) at (26.5,-62) {病害\\$\hat y_2$};
  \node[output] (outa3) at (44.5,-62) {定位\\$\hat y_3$};
  \draw[wire] (shared.south) -- (26.5,-31);
  \draw[flow] (26.5,-31) -| (head1.north);
  \draw[flow] (26.5,-31) -- (head2.north);
  \draw[flow] (26.5,-31) -| (head3.north);
  \draw[flow] (head1.south) -- (outa1.north);
  \draw[flow] (head2.south) -- (outa2.north);
  \draw[flow] (head3.south) -- (outa3.north);

  \node[prior] (prior) at (84.5,-20)
    {共同先验\\$\bm\mu,\bm\Omega$};
  \node[model] (modelb1) at (66.5,-43) {$f_{\bm w_1}$};
  \node[model] (modelb2) at (84.5,-43) {$f_{\bm w_2}$};
  \node[model] (modelb3) at (102.5,-43) {$f_{\bm w_3}$};
  \node[output] (outb1) at (66.5,-62) {类别\\$\hat y_1$};
  \node[output] (outb2) at (84.5,-62) {病害\\$\hat y_2$};
  \node[output] (outb3) at (102.5,-62) {定位\\$\hat y_3$};
  \draw[stat] (prior.south) -- (84.5,-31);
  \draw[stat] (84.5,-31) -| (modelb1.north);
  \draw[stat] (84.5,-31) -- (modelb2.north);
  \draw[stat] (84.5,-31) -| (modelb3.north);
  \draw[flow] (modelb1.south) -- (outb1.north);
  \draw[flow] (modelb2.south) -- (outb2.north);
  \draw[flow] (modelb3.south) -- (outb3.north);

  \node[model] (modelc1) at (124.5,-20) {$f_1$};
  \node[model] (modelc2) at (142.5,-20) {$f_2$};
  \node[model] (modelc3) at (160.5,-20) {$f_3$};
  \node[output] (outc1) at (124.5,-43) {类别\\$\hat y_1$};
  \node[output] (outc2) at (142.5,-43) {病害\\$\hat y_2$};
  \node[output] (outc3) at (160.5,-43) {定位\\$\hat y_3$};
  \node[relation] (relation) at (142.5,-63)
    {同对象预测输出约束\\$L_{\mathrm{rel}}$};
  \draw[flow] (modelc1.south) -- (outc1.north);
  \draw[flow] (modelc2.south) -- (outc2.north);
  \draw[flow] (modelc3.south) -- (outc3.north);
  \draw[constraint] (outc1.south) -- (124.5,-53) -- ($(relation.north)+(-10,0)$);
  \draw[constraint] (outc2.south) -- (relation.north);
  \draw[constraint] (outc3.south) -- (160.5,-53) -- ($(relation.north)+(10,0)$);

  \node[heading] at (0,-85) {(d) 表示/模型参数 × 条件前向：专家由各任务共同训练};
  \node[prior,text width=16mm] (input) at (12,-116)
    {同一图像\\$\bm x$};
  \node[model,text width=25mm] (expert1) at (54,-101)
    {专家1：$\bm e_1(\bm x)$};
  \node[model,text width=25mm] (expert2) at (54,-116)
    {专家2：$\bm e_2(\bm x)$};
  \node[model,text width=25mm] (expert3) at (54,-131)
    {专家3：$\bm e_3(\bm x)$};
  \node[prior,text width=28mm] (gate) at (54,-151)
    {任务$k$的门控\\$\bm g_k(\bm x)$};
  \node[model,text width=24mm,minimum height=16mm] (mix) at (111,-116)
    {按权重混合\\$\bm z_k$};
  \node[model,text width=24mm,minimum height=16mm] (task) at (151,-116)
    {任务头$q_k$\\预测输出$\hat y_k$};
  \node[output,text width=30mm] at (151,-134)
    {$k=1,2,3$\\类别、病害、定位};
  \draw[wire] (input.east) -- (28,-116);
  \draw[flow] (28,-116) |- (expert1.west);
  \draw[flow] (28,-116) -- (expert2.west);
  \draw[flow] (28,-116) |- (expert3.west);
  \draw[flow] (28,-116) |- (gate.west);
  \draw[flow] (expert1.east) -- (mix.north west);
  \draw[flow] (expert2.east) -- (mix.west);
  \draw[flow] (expert3.east) -- (mix.south west);
  \draw[flow] (gate.east) -- (111,-151) -- (mix.south);
  \node[anchor=west] at (114,-145) {混合权重};
  \draw[flow] (mix.east) -- (task.west);
\end{tikzpicture}

  \caption{共同计算、模型参数先验与预测输出约束的共享位置。上排固定类别、病害和定位任务，省略重复预测输入；实线箭头是前向计算，虚线是先验联系或训练约束。先验面板只展开同维同坐标模型参数块$\bm w_k$的联系，未画异构预测输出头。预测输出面板示意训练约束，并未画出预测传递所需的额外前向路径。下排将共同计算展开为条件混合：同一图像送入专家与任务$k$的门控，再由任务权重混合专家表示；此处$\bm g_k(\bm x)$是由$g_{k,m}$组成的门控向量，不是共享梯度。三个任务分别设置门控和头，损失反传到对应部件及共同专家；反向路径未画。稠密门控仍需计算全部专家。}
  \label{fig:ket-sharing-locations}
\end{figure*}
\FloatBarrier

\section{联合学习目标的协调}
\label{sec:ket-sharing-objectives}

选好分组、搭好共享结构之后，一个绕不开的操作问题出现了：三个任务的损失都回传到同一骨干，这一步到底该怎么更新？这就是联合目标的协调问题。为了把问题看清楚，先暂时固定任务分组、网络结构和采样方式，只聚焦三件互相独立的事：损失的数值尺度、梯度的大小、梯度的方向。它们绝不能互相替代，这里有三个常见错觉要先打掉：一个损失数值很大，未必说明它对应的参数正在快速变化；一个梯度范数很大，未必说明那个任务更重要；而此刻两个梯度方向冲突，只反映这一步更新的局部矛盾，还不能据此认定长期存在负迁移。

技术上有两条互补的路线：一是先给各项损失配上权重、加成一个标量目标，再对它求梯度（加权求和派）；二是保留每项任务各自的梯度，直接在梯度层面构造一个联合更新方向（梯度协调派）。二者并非互斥——梯度方法解到最后往往也能等价地看成某种加权。事实上，如果系数相同且在同一次求导中被当作常量，“对加权损失求梯度”和“对梯度加权求和”得到的是同一个向量。所以真正值得比较的不是表面形式，而是：这些系数怎么确定、多久重新估计一次、以及得到的更新方向满足什么性质（是只优化平均，还是保证谁都不退化）。

\subsection{任务权重的确定}
\label{sec:ket-sharing-weights}

最直接的协调手段就是给损失加权。假设当前批次恰好包含所有任务的监督，用非负权重 $w_k$ 表示每项损失对这次更新的贡献大小，加权联合目标写成：
\begin{equation}
  L_{\mathrm{joint}}
  =\sum_{k=1}^{K}w_kL_k,
  \qquad w_k\geq0.
  \label{eq:ket-sharing-weighted-loss}
\end{equation}
权重固定时，共享梯度就是各任务梯度的加权和 $\sum_k w_k\bm g_k$。但调权重前必须先解决一个更基础的问题——损失尺度，否则权重根本无从解释。放大某项的权重可能出于两种完全不同的动机：表达业务优先级（这个任务更重要），或仅仅补偿数值尺度（它的单位导致数值偏大/偏小），这两者必须分别记录、不能混为一谈。一个极有说服力的例子：同样是预测一个物理位置，把坐标单位从米换成厘米，平方误差的数值和它的梯度会瞬间放大 $10^4$ 倍（因为 1 米 = 100 厘米，平方后是一万倍）；可类别交叉熵毫无变化，定位任务的真实责任也没有任何改变。如果不先归一化，模型会误以为定位“重要一万倍”。所以正确顺序是：先统一单位、讲清归一化方式、明确损失是按样本还是按像素平均，然后才谈得上用权重表达取舍。

权重还应和前面确定的交付责任一一对应：三项共同交付时，权重设置要服从“各项都不能跌破最低标准”的约束，不能让平均好看；只部署类别时，辅助任务权重的大小就该由类别任务的验证结果来决定；若是多个地区副本，则要说清是各地区等权、还是按各自的服务请求量加权。固定权重的好处是给出一个明确、可复现的优化目标，但它在训练的不同阶段未必一直合适（比如初期和后期任务的难度关系会变）。动态调整权重能跟随训练状态变化，听起来更聪明，但它只能依据可计算的训练信号调整，绝不能越权替你决定业务上谁更重要——优先级必须由人给定。

怎么评价一个“自动调权重”的方法到底有没有价值？关键是配一个足够强又公平的基线：在完全相同的训练样本、网络结构和基础优化器下，和简单的等权联合训练比较；并在预先约定的搜索预算内去调固定权重或采样比例、保持停止规则一致，才能看出专用协调器究竟额外带来了什么。注意单任务基线在这里并不合适用来隔离“调权”的作用，因为它连共享结构都一起变了。成本也要算全：搜索权重、为每个任务单独反传梯度、求解额外子问题的开销都得记账。至于“改变采样比例”其实是从另一个角度改变目标，这一点见第~\ref{sec:ket-sharing-sampling}节。

第一个自动调权的直觉是“帮一把进步慢的学生”：如果某个任务的损失下降得比别的任务慢，就该在近期多给它一点权重。\term{动态权重平均}（Dynamic Weight Average, DWA）把这个想法落地为：比较任务损失在相邻两个训练阶段的比值，用它生成当前权重。为避免被单批次噪声主导，$L_k^{(t)}$ 取的是一个阶段内的平均损失、而不是某个随机小批次的值，并要求分母为正。具体令
\[
  r_k^{(t-1)}=\frac{L_k^{(t-1)}}{L_k^{(t-2)}},
  \qquad
  w_k^{(t)}
  =K\frac{\exp(r_k^{(t-1)}/T)}
          {\sum_j\exp(r_j^{(t-1)}/T)},
\]
其中温度 $T>0$ 控制权重之间拉得有多开（$T$ 越大越趋平均），归一化保证权重之和为 $K$（与等权时总量一致）；还没有两个历史阶段可用时，就用预先设定的初值。它的逻辑是损失下降越慢、$r_k$ 越大、分到的权重越高。但这个依据有几个失真来源：尺度突变、阶段内残余的批次噪声、以及接近零的分母都可能让比值剧烈抖动。更要认清 DWA 的边界：它只在平衡“相对下降速度”，既不判断任务本身重不重要，也完全不处理梯度方向冲突——两个任务哪怕进步一样快，梯度仍可能正相反\scite{ket-liu2019-mtan}。

第二种依据更有概率味道：不再凭进度猜权重，而是把权重和“这个任务的标签本身有多噪”联系起来。\term{不确定性加权}（uncertainty weighting）的核心思想是，如果一个任务的观测天然带有很大噪声，那就不该要求模型去死拟合那些噪声，理应给它较小的有效权重；这个噪声尺度应该能从数据里学出来。先在回归任务上推导：取 $d_k$ 维连续输出，假设任务预测 $f_k(\bm x)$ 与真实目标之间的残差是独立高斯噪声：
\begin{equation}
  y_k=f_k(\bm x)+\bm\varepsilon_k,
  \qquad
  \bm\varepsilon_k\sim
  \mathcal N(\bm0,\sigma_k^2\bm I).
\end{equation}
这里 $\sigma_k>0$ 对同一个任务的所有输入都相同，称为“任务级同方差噪声尺度”——它刻画的是这个任务整体的标注可靠程度，绝不是给每张图打的置信度分数（这点特别容易误解）。把逐样本平方残差范数的一半在批次上平均记作 $L_k^{\mathrm{reg}}$，并令 $s_k=\log\sigma_k^2$（用对数参数化保证方差为正）。把高斯负对数密度展开后，单任务平均负对数似然变成一个关于模型和 $s_k$ 的目标：
\begin{equation}
  \mathcal L_k^{\mathrm{reg}}
  =e^{-s_k}L_k^{\mathrm{reg}}
   +\frac{d_k}{2}s_k+C.
  \label{eq:ket-sharing-uncertainty-regression}
\end{equation}
这个式子两项的博弈正是方法的精髓：第一项 $e^{-s_k}L_k^{\mathrm{reg}}$ 按逆方差给残差加权——任务越噪（$s_k$ 越大），拟合残差的压力越小；第二项 $\frac{d_k}{2}s_k$ 则是个“刹车”，防止模型走极端：如果只靠把方差调到无穷大就能让第一项归零、从而彻底忽略这个任务，那还要第二项做什么？正是它随方差增大而增大，逼模型在“拟合数据”和“声称噪声大”之间诚实权衡。$d_k/2$ 来自 $d_k$ 个输出维的高斯归一化因子；若你改成按维度平均，整个似然项的尺度要相应调整。此外，对每个任务分别取样本平均也是一种归一化选择，当各任务样本数不同时，这并不等于直接对全部观测的似然求和，两者隐含的权重不同。

分类任务的输出是类别概率、归一化方式和回归不同，不能直接套用上面的回归式。记 $\bm f_k$ 为未归一化的类别分数（logits），$f_{k,c}$ 为类别 $c$ 那个分量，$L_k^{\mathrm{CE}}$ 为原始分数上的交叉熵。把分类概率改成带温度的 $\operatorname{softmax}(e^{-s_k}\bm f_k)$（即给 logits 乘上一个可学的逆温度），展开单样本负对数概率，得到：
\begin{align}
  \mathcal L_k^{\mathrm{cls}}
  ={}&e^{-s_k}L_k^{\mathrm{CE}}
  +\log\sum_c\exp(e^{-s_k}f_{k,c})
  \nonumber\\
  &-e^{-s_k}\log\sum_c\exp(f_{k,c}),
  \label{eq:ket-sharing-uncertainty-classification}
\end{align}
后两个对数项完全来自 softmax 的配分函数（归一化分母），不能为了简洁随意删掉。Kendall 等人在实践中对这个难处理的归一化因子做了近似，得到广泛使用的简洁形式 $e^{-s_k}L_k^{\mathrm{CE}}+\frac12s_k$（结构和回归式很像）；但必须知道这个简洁式是近似产物，并不是对任意类别分数都成立的恒等式\scite{kendall2018}。所以写论文或报告时应当说明：你用的是精确的分类概率似然，还是那个被近似过的常用目标。

实现上，选定似然形式后可以从 $s_k=0$（即初始方差为 1）开始，每个批次同时更新模型参数和这些噪声参数。为了看清新加入的 $s_k$ 到底学到了什么，把单任务目标抽象成 $e^{-s_k}L_k+c_ks_k$（$c_k>0$ 是对应的对数项系数）。固定模型参数、只对 $s_k$ 求导：
\begin{equation}
  \frac{\partial\mathcal L_k}{\partial s_k}
  =-e^{-s_k}L_k+c_k.
\end{equation}
令导数为零，在 $L_k>0$ 时解出 $\sigma_k^2=L_k/c_k$：学出的噪声方差恰好正比于当前的误差尺度。这就把它的身份讲透了——不确定性加权本质是在自动匹配每个任务的误差量级，它并不是在独立表达“哪个任务更重要”。由此也带来风险：当某任务误差趋于零时，方差可能被推向零边界、权重趋于无穷，需要加数值下界或先验稳住；还得回头检查高斯噪声假设是否合理、有没有“其实重要但困难”的任务被过早降权放弃，以及任务级平均是否掩盖了某些稀缺群体上的大误差。

就算把损失尺度归一化好了，各任务回传的梯度范数仍可能差很多——一个任务梯度巨大、主导每次更新，另一个几乎没声音。梯度归一化（Gradient Normalization, GradNorm）同时利用“梯度大小”和“相对学习进度”两个信号来调权，目标是让下降偏慢的任务在共享层获得更大的实际梯度作用。它需要先选定一个测量块 $\bm{\theta}_{\mathrm b}$（通常取各任务分叉之前的共享层），并记任务 $k$ 在该块上、尚未乘权重的原始梯度为 $\bm g_{k,\mathrm b}^{(t)}$。将当前权重视为正的常量，测得
\begin{equation}
  G_k^{(t)}
  =\left\|\nabla_{\bm{\theta}_{\mathrm b}}
     \bigl(w_k^{(t)}L_k^{(t)}\bigr)\right\|_2
  =w_k^{(t)}\|\bm g_{k,\mathrm b}^{(t)}\|_2,
  \label{eq:ket-sharing-gradnorm-size}
\end{equation}
这里有个实现细节值得点出：必须保证权重为正。一个干净的做法是用 $w_k=\exp(v_k)$ 参数化（$v_k$ 可在整个实数域自由更新，指数天然为正），或者在每次更新后把权重投影回正数域；不能在做了无约束更新后还想当然地认为正性自动保持。接下来要区分“绝对损失大”和“下降得慢”这两件不同的事：先用非零的初始损失 $L_k^{(0)}$ 把当前损失归一化，再拿它和所有任务的平均进度相比：
\begin{equation}
  q_k^{(t)}=\frac{L_k^{(t)}}{L_k^{(0)}},
  \qquad
  r_k^{(t)}
  =\frac{q_k^{(t)}}{K^{-1}\sum_jq_j^{(t)}}.
  \label{eq:ket-sharing-gradnorm-rate}
\end{equation}
其中 $q_k$ 是当前损失占自身初值的比例（越小说明相对进步越大），$r_k$ 是它相对所有任务平均进度的比值。关键解读是：$r_k>1$ 只意味着这个任务“比平均水平下降得慢”，并不代表它的绝对损失更大——一个原本很容易、绝对损失很小的任务，若迟迟没继续下降，照样会有高 $r_k$。这也带来边界情况：初值为零或极小时要预先另设有效参照（否则除以近乎零），若此刻所有任务损失都为零，平均进度为零，本步就该停止调权，而不是硬去除以零。

记平均加权范数为$\bar G=K^{-1}\sum_kG_k$。以$\bar G r_k^\alpha$作为任务$k$的目标范数，$\alpha\geq0$决定对慢任务倾斜多少；让测量范数接近这个目标，可最小化
\begin{equation}
  L_{\mathrm{grad}}
  =\sum_{k=1}^{K}
  \left|G_k-
    \operatorname{stopgrad}(\bar G r_k^\alpha)\right|.
  \label{eq:ket-sharing-gradnorm-objective}
\end{equation}
这个辅助损失用绝对值衡量“实测梯度范数”和“目标范数”的差距。实现时有两个“停梯度”纪律：目标值 $\bar G r_k^\alpha$ 里的量在本次调权时要当常量（$\operatorname{stopgrad}$），各任务已算出的梯度本身也固定，只对权重 $w_k$ 求导；否则若让目标随权重一起变化、或者干脆拿这个辅助损失去训练表示，就彻底改变了各方的职责、得到错误的东西。模型主体仍按加权后的任务损失更新，而权重在每步调完后要保持为正并归一化到和为 $K$。通过这套机制，GradNorm 把“梯度大小”和“相对学习速度”两个信号都用上了\scite{ket-chen2018-gradnorm}。

用一组数字感受它：设三个任务的相对进度 $q_k=(0.2,0.5,0.8)$，平均 0.5，于是 $r_k=(0.4,1,1.6)$——第三个任务（定位）下降最慢。再设平均梯度范数 $\bar G=2$、$\alpha=1$，则三个任务被推向的目标范数为 $(0.8,2,3.2)$，定位确实拿到最大的目标梯度。但请记住这只是在调节“训练努力的分配”，绝不等于证明“定位在业务上应该享有更高优先级”。如果定位下降慢纯粹是因为它的标签噪声大，那么持续给它加力只是在浪费算力去拟合噪声。测量层选在哪、初值怎么取、批次波动都会影响权重。最后，不确定性加权依据的是“误差噪声尺度”、GradNorm 依据的是“梯度范数和进度”，两者机理不同，无论用哪个都要回到逐任务的独立验证来检查真实效果。

GradNorm 这类方法有个实际成本：要为每个任务单独反传、取得它在共享层的梯度范数，当任务很多时这笔额外反传很贵。快速自适应多任务优化（Fast Adaptive Multitask Optimization, FAMO）换了个更省的依据：它不显式求梯度，而是用任务损失（取正值）的“相对改善量”来组织协调，并把权重这个子问题的求解分摊到连续的多个训练步上，避免每步解一个完整优化。它利用一个小步长下的近似：相邻两个模型状态之间的对数损失变化 $\log L_k(\bm\theta^{(t)})-\log L_k(\bm\theta^{(t+1)})$，可以近似替代调权时本来需要的“梯度内积反馈”，据此逐步调整非负且归一化的权重；模型主体仍按构造出的加权总损失正常反传。这样就不必每步为每个任务保存一整份共享梯度，但它仍需拿到各任务损失、维护一组任务标量权重、并跑各任务的预测分支，所以并非零成本。使用时还要固定若干条件：损失必须为正并设数值下界、前后两次损失要在相同条件下测量、反馈近似才成立；既不能把总开销写成与任务数无关，也不能据此承诺“任意有限步内每个任务损失都下降”\scite{ket-liu2023-famo}

\subsection{联合更新方向的构造}
\label{sec:ket-sharing-directions}

GradNorm 把各任务梯度的大小调到相近之后，问题并没有完全解决：大小相当的两个梯度，方向仍可能正好相反，合成更新还是会厚此薄彼。这一小节就专门处理“方向”。下面固定私有参数、省略阶段上标，用 $\bm d$ 表示由各任务梯度合成出的共享更新方向，实际参数更新为 $\bm{\theta}_{\mathrm s}^{+}=\bm{\theta}_{\mathrm s}-\eta\bm d$（$\eta>0$）。判断这一步到底帮了谁、害了谁，只需在梯度局部 Lipschitz 的假设下，把任务损失沿该方向做一阶展开：
\begin{equation}
  L_k(\bm{\theta}_{\mathrm s}-\eta\bm d)
  =L_k(\bm{\theta}_{\mathrm s})
   -\eta\bm g_k^{\top}\bm d+O(\eta^2).
  \label{eq:ket-sharing-direction-effect}
\end{equation}
主导项是 $-\eta\bm g_k^\top\bm d$：只要内积 $\bm g_k^\top\bm d$ 为正，足够小的一步就一定降低任务 $k$ 当前的损失；为负则会抬高它。这个简单判据给出了三种由弱到强的处理思路，正好对应接下来三个小节：一是“消冲突”，把每个梯度里妨碍别人下降的分量删掉（PCGrad）；二是“找共赢”，直接构造一个让所有任务内积同时为正的共同下降方向（MGDA）；三是“谈分配”，在众多共同下降方向里进一步决定谁多分一点收益（Nash 协商）。这三种要求可以叠加，并不是三个互斥又穷尽的格子。

\subsubsection{基于冲突消解的方向构造}
\label{sec:ket-sharing-pcgrad}

第一种思路很外科手术：既然任务 $k$ 的梯度里有一部分正顶着任务 $j$（两者内积为负），那在把它们相加之前，先把 $k$ 梯度中“妨碍 $j$ 下降”的那个分量切掉。投影冲突梯度（projecting conflicting gradients, PCGrad）的做法就是把 $k$ 梯度的一个副本投影到 $j$ 梯度的正交平面（垂直于 $j$ 的方向）上，使投影后两者不再冲突。具体为任务 $k$ 复制一份原始共享梯度 $\widetilde{\bm g}_k=\bm g_k$；逐个考察其他任务 $j$，用当前副本与“原始的” $\bm g_j$ 的内积判断是否冲突。若内积为负且 $\bm g_j$ 非零，就执行投影：
\begin{equation}
  \widetilde{\bm g}_k\leftarrow
  \widetilde{\bm g}_k-
  \frac{\widetilde{\bm g}_k^{\top}\bm g_j}
       {\|\bm g_j\|_2^2}\bm g_j.
  \label{eq:ket-sharing-pcgrad-projection}
\end{equation}
减去的那一项，正是当前副本在 $\bm g_j$ 方向上的投影分量；去掉它之后，副本与 $\bm g_j$ 的内积恰好为零，即不再有正对着 $\bm g_j$ 的冲突成分。没有冲突（内积非负）时就保持原样，遇到零梯度则跳过以免除零。每个任务的副本都按某个随机顺序逐一对照其他任务做投影，最后把处理过的各副本求和或按约定平均；私有参数则依旧只用本任务自己的监督更新。一个实现要点：投影时参照的永远是其他任务的“原始”梯度，而不是它们已经被投影改过的副本，否则结果会随处理顺序纠缠不清\scite{ket-yu2020-pcgrad}。

接下来用同一组二维二次损失作为统一的“教学沙盘”，把投影、共同下降、协商三种方向构造都亲手算一遍、放在一起对比。设 $\bm\theta=(\theta_1,\theta_2)^\top$ 是两个共享自由度，两个任务各自希望参数靠近不同的中心点：
\begin{align}
  L_1(\bm\theta)
    &=\tfrac12\|\bm\theta-(-1,0)^{\top}\|_2^2,
  \nonumber\\
  L_2(\bm\theta)
    &=\tfrac12\|\bm\theta-(2,-2)^{\top}\|_2^2.
  \label{eq:ket-sharing-toy-losses}
\end{align}
再次声明这些二次损失只是用来模拟共享块上的竞争，是教学构造、不是真实花卉测量。在原点处，两个任务的梯度分别为 $\bm g_1=(1,0)^\top$、$\bm g_2=(-2,2)^\top$，内积为 $-2<0$，确实冲突。先看最朴素的直接求和：$\bm d_{\mathrm{sum}}=(-1,2)^\top$。它和 $\bm g_1$ 的内积是 $-1<0$，意味着沿这个“总和下降”的方向走一小步，任务 1 的损失反而会上升——总损失可能下降，但这是靠牺牲任务 1 换来的，这正是需要协调的症结。

将这两个梯度代入式~\eqref{eq:ket-sharing-pcgrad-projection}，各自删去相对另一任务的冲突分量：
\begin{align}
  \widetilde{\bm g}_1
    &=(1,0)^{\top}-\frac{-2}{8}(-2,2)^{\top}
      =(0.5,0.5)^{\top},
  \nonumber\\
  \widetilde{\bm g}_2
    &=(-2,2)^{\top}-\frac{-2}{1}(1,0)^{\top}
      =(0,2)^{\top}.
\end{align}
把两个投影后的向量相加，得到 $\bm d_{\mathrm{PC}}=(0.5,2.5)^\top$。验证它和两个原始梯度的内积分别为 $0.5$ 和 $4$，都为正，说明这一步对两个任务都是一阶下降。光看一阶符号还不够，再固定步长 $\eta=0.1$，把更新后的参数真实代回二次损失（而不是只停留在一阶近似），得到更新前后的损失对比：
\begin{align}
  (L_1,L_2)(\bm0)&=(0.5,4),
  \nonumber\\
  (L_1,L_2)(-\eta\bm d_{\mathrm{sum}})
    &=(0.625,3.425),
  \nonumber\\
  (L_1,L_2)(-\eta\bm d_{\mathrm{PC}})
    &=(0.4825,3.6325).
  \label{eq:ket-sharing-pcgrad-numbers}
\end{align}
结果很说明问题：直接求和把任务 1 的损失从 $0.5$ 抬高到 $0.625$（损人利己）；而 PCGrad 投影后任务 1 降到 $0.4825$、任务 2 也降到 $3.6325$，两项都下降。当然要克制解读：这里比较的只是固定步长下的当前损失，而且两个方向向量长度不同，不能据此排出一个一般性的“最终预测性能”名次。

PCGrad 也不是没有短板。逐对投影存在顺序依赖：先对 $j$ 投影、再对 $l$ 投影时，后一次投影可能又改变了前面好不容易调好的与 $j$ 的夹角，所以它的结果并不是“同时满足所有正交约束”的精确解。批次里出现的负内积还可能只是梯度估计噪声，未必是真冲突。此外，为每个任务分别算梯度、还要逐对做内积和投影，训练开销明显增加；而且它改善的只是“当前这一步”的更新行为，并不保证训练结束后所有任务都能超过同等预算的单任务基线。

\subsubsection{基于共同下降的方向构造}
\label{sec:ket-sharing-mgda}

PCGrad 是成对地、逐个地消解冲突；如果想一次性保证“这一步同时降低所有任务的损失”，可以用更有原则的几何构造。\term{多梯度下降算法}（multiple gradient descent algorithm, MGDA）的候选方向是所有任务梯度的凸组合——用一组非负、和为 1 的系数把各梯度混起来；在这些组合里，它选择长度最短（离零点最近）的那一个。背后的直觉是：凸包里最靠近原点的那个向量，往往正好指向“谁都不得罪”的共赢方向。形式化为：
\begin{align}
  \bm\alpha^\star
  &=\argmin_{\substack{\alpha_k\geq0\\\sum_k\alpha_k=1}}
  \frac12\left\|\sum_{k=1}^{K}\alpha_k\bm g_k\right\|_2^2,
  \nonumber\\
  \bm d^\star&=\sum_k\alpha_k^\star\bm g_k.
  \label{eq:ket-sharing-mgda}
\end{align}
目标里的平方范数只是把“离原点最近”写成一个标准、可快速求解的凸二次最小二乘问题（在系数单纯形上）。一定要区分这里的 $\alpha_k^\star$ 和前面讲的业务权重：它是为了寻找共同下降方向而解出来的几何系数，完全不是人预先指定的任务重要性——某任务系数为零，也不代表它不重要，只说明当前最优共赢方向不需要它那一分量。Sener 和 Koltun 把这一构造用到共享网络的多目标更新中，并研究了降低其求解成本的近似办法\scite{ket-sener2018-mgda}。

为什么“凸包里离原点最近的点”就一定是共同下降方向？给一个简短的几何论证：从最优点 $\bm d^\star$ 朝凸包的任意一个顶点 $\bm g_k$ 移动，得到的点仍在凸包内；由于 $\bm d^\star$ 已是最近点，朝任何顶点移动都不能让范数更小，所以这个方向上的单侧方向导数必须非负。对平方范数 $\frac12\|\bm d\|^2$ 写出这个最优性条件：
\begin{equation}
  (\bm g_k-\bm d^\star)^{\top}\bm d^\star\geq0,
  \qquad
  \bm g_k^{\top}\bm d^\star
  \geq\|\bm d^\star\|_2^2.
  \label{eq:ket-sharing-common-descent}
\end{equation}
当 $\bm d^\star
\ne\bm0$ 时，右侧 $\|\bm d^\star\|_2^2$ 严格为正，于是对所有任务都有 $\bm g_k^\top\bm d^\star>0$，结合前面的一阶判据，所有任务沿 $-\bm d^\star$ 走足够小一步都下降——这就是“共同下降”的来历。不过这只是一阶结论，要保证有限步长后真的下降，还得管住二阶余项；若梯度在邻域内有光滑常数 $\beta_k>0$，标准的下降引理给出上界：
\begin{equation}
  L_k(\bm\theta-\eta\bm d)
  \leq L_k(\bm\theta)
  -\eta\bm g_k^{\top}\bm d
  +\frac{\beta_k\eta^2}{2}\|\bm d\|_2^2.
\end{equation}
其中 $\bm\theta$ 是当前共享块。把共同下降方向 $\bm d^\star$ 代入，只要步长满足 $0<\eta<2/\max_k\beta_k$ 且更新后仍停在该光滑邻域内，二次上界就严格小于原损失，得到真正的有限步共同下降。要注意这个结论不随大步长迁移——步子迈太大，二阶项反扑，照样可能损害某任务。相关的光滑性和凸优化基础可参考\BookChapterRef{chap:convex-learning-and-stability}{卷3《基础、理论与可信性》的“一般学习的可学习性”章}。

回到式~\eqref{eq:ket-sharing-toy-losses}的零点，以$\alpha_1=a,\alpha_2=1-a$表示线段，组合为$(3a-2,2-2a)^{\top}$。在$0\leq a\leq1$上最小化平方范数，导数为零得到$a=10/13$，从而
\begin{equation}
  \bm d^\star=(4/13,6/13)^{\top},
  \qquad
  \bm g_1^{\top}\bm d^\star
  =\bm g_2^{\top}\bm d^\star
  =4/13.
  \label{eq:ket-sharing-mgda-numbers}
\end{equation}
两个内积都等于 $4/13>0$，验证了该方向在原点处确实是共同下降方向。但 MGDA 也有走不动的时候：把参数移到两个任务中心的中点 $\bm\theta=(0.5,-1)^\top$ 后，两个任务的梯度变成 $(1.5,-1)^\top$ 和它的相反向量（等大反向），它们的等权凸组合恰为零向量。此时无论选什么方向，它与这两个梯度的内积都互为相反数，不可能同时严格为正——共赢的一阶方向不复存在。

这就引出多目标优化的两个核心概念。在无约束可微情形下，当零向量落在各任务梯度的凸包内（凸组合能得到零）时，称当前点为\term{帕累托驻点}（Pareto stationary point）：它的含义是不存在任何一阶方向能让所有目标同时严格下降。而\term{帕累托有效}（Pareto efficient）是一个更强的、关于实际可行解的全局概念：不存在另一个解能让所有目标都不差、且至少一个目标严格更好。两者不能混：驻点只是局部的一阶必要条件，并不保证全局有效；在非凸神经网络里，优化完全可能停在一个很差的帕累托驻点。另外要注意作用范围：上面的判断只针对固定私有参数时的共享块（任务头还能继续更新），而且无论是驻点还是有效，都不等于“已经超过了单任务基线”——那是另一码事。

求精确组合昂贵时，可在系数单纯形上用Frank--Wolfe迭代。当前$\bm d=\sum_k\alpha_k\bm g_k$，先按线性化目标选择内积$\bm g_j^{\top}\bm d$最小的顶点$\bm e_j$，再在$0\leq\gamma\leq1$上搜索并更新$\bm\alpha\leftarrow(1-\gamma)\bm\alpha+\gamma\bm e_j$。用当前范数与最小顶点内积之差衡量求解误差：
\begin{equation}
  \Delta=\|\bm d\|_2^2-\min_k\bm g_k^{\top}\bm d.
\end{equation}
若$\Delta<\|\bm d\|_2^2$，则最小顶点内积也为正，当前近似方向仍共同一阶下降。只设固定迭代次数而不检查误差，不能沿用精确解的判据。

另一种节省方式是减少所保存梯度的维数。MGDA-UB在共享表示上求一个上界问题：设$\bm z=h_{\bm{\theta}_{\mathrm s}}(\bm x)$，$\bm r_k=\nabla_{\bm z}L_k$，$\bm J=\partial\bm z/\partial\bm{\theta}_{\mathrm s}$为雅可比，则链式法则给出$\bm g_k=\bm J^{\top}\bm r_k$。以相同凸系数构造模型参数方向时，真实下降判据中的内积为
\begin{equation}
  \bm g_k^{\top}\bm d
  =\bm r_k^{\top}\bm J\bm J^{\top}
    \sum_j\alpha_j\bm r_j.
  \label{eq:ket-sharing-mgda-feature-metric}
\end{equation}
因此，直接最小化表示梯度范数并不是同一个度量。其计算依据是谱范数对任意$\bm r(\bm\alpha)=\sum_k\alpha_k\bm r_k$给出的上界：
\[
  \left\|\sum_k\alpha_k\bm g_k\right\|_2
  =\|\bm J^\top\bm r(\bm\alpha)\|_2
  \leq\|\bm J\|_2\,\|\bm r(\bm\alpha)\|_2.
\]
$\|\bm J\|_2$在同一模型参数状态和批次下不随组合系数改变，故可通过表示范数控制模型参数范数的上界。但上界最小不保证真实范数最小，更不直接保证各任务真实内积正。

为单独检查度量变化，另作二维教学构造：这次把$\bm r_1=(1,0)^{\top}$、$\bm r_2=(-2,2)^{\top}$作为表示梯度，在表示空间求得的最短组合系数为$(10/13,3/13)$。若对角雅可比为$\bm J=\operatorname{diag}(10,1)$，它把第一坐标的梯度放大10倍，模型参数梯度与同系数组合变成
\[
  \begin{aligned}
    \bm g_1&=(10,0)^{\top},\qquad
    \bm g_2=(-20,2)^{\top},\\
    \bm d&=(40/13,6/13)^{\top}.
  \end{aligned}
\]
此时$\bm g_2^{\top}\bm d=-788/13<0$，沿$-\bm d$的小步反而提高任务2的损失。虽然$\bm J$满秩，两个坐标的相对尺度已经改变，表示空间的下降判断不能直接搬到模型参数空间。

若$\bm J\bm J^{\top}$在这些梯度张成的子空间上等于$c\bm I$且$c>0$，两种度量只差正常数，可继承下降关系；一般情况须检查方法假设或复查真实内积。即使真实批次内积全正，结论也只针对当前批次，不承诺总体风险逐步降低。

\subsubsection{基于协商收益的方向构造}
\label{sec:ket-sharing-nash}

MGDA 解决了“找一个共赢方向”，但共同下降方向通常有无穷多个，同样是共赢，有的对任务 1 好处多、有的对任务 2 好处多，还需要一个原则来决定“一阶收益怎么在任务间分配”。Nash-MTL 借用博弈论里的协商（bargaining）思想：以“完全不更新”作为各方的参考状态，让各任务按各自能拿到的一阶效用来协商\scite{ket-navon2022-nash}。为了只比较方向、不被向量长度任意放大影响，固定 $\|\bm d\|_2\leq1$、省去共同步长，把任务 $k$ 的收益定义为它的一阶下降量 $u_k(\bm d)=\bm g_k^\top\bm d$。它只接受所有任务收益都为正的方向（先保证共赢），然后最大化各任务收益的对数之和：
\begin{equation}
  \max_{\substack{\|\bm d\|_2\leq1\\
                  \bm g_k^{\top}\bm d>0\ \forall k}}
  \sum_{k=1}^{K}\log(\bm g_k^{\top}\bm d).
  \label{eq:ket-sharing-nash-objective}
\end{equation}
为什么用对数和而不是普通求和？因为 $\sum_k\log u_k$ 等价于最大化收益的乘积 $\prod_k u_k$，而乘积目标有一个公平性特质：它绝不允许牺牲者——若某项收益趋近零，整个目标就趋向负无穷，所以你无法用“让某任务收益为负”去换其他任务收益更大。这正契合协商的精神：每一方相对“不更新”都必须受益。要认清它比较的是给定损失下的局部线性收益，不是直接在比最终的验证分数。还有个可行性前提：如果当前根本不存在严格的共同下降方向（比如已在帕累托驻点），可行集为空，就不该硬套这个协商目标。

把共享梯度列成$\bm G=[\bm g_1,\ldots,\bm g_K]$，其Gram矩阵$\bm H=\bm G^{\top}\bm G$记录任务两两内积。存在严格可行方向时，向外缩放增加全部效用，所以最优解位于单位球面。用$\lambda$作球约束乘子，关于方向求导有
\begin{equation}
  \sum_k\frac{\bm g_k}{\bm g_k^{\top}\bm d}
  =\lambda\bm d.
\end{equation}
与单位方向$\bm d$取内积后，左侧各项都是1，得到$\lambda=K$。再定义正系数$a_k=1/(\sqrt K\,\bm g_k^{\top}\bm d)$，驻点关系写成$\bm d=\bm G\bm a/\sqrt K$。将它代回效用，得到系数必须满足的方程：
\begin{equation}
  \bm H\bm a=\bm a^{-1},
  \qquad
  (\bm a^{-1})_k=1/a_k.
  \label{eq:ket-sharing-nash-system}
\end{equation}
推导给出一个干净的结构：在严格可行时，最优解一定在单位球面上（因为放大方向能同时提高所有收益），最后可写成 $\bm d_{\mathrm N}=\bm G\bm a/\sqrt K$，其中正系数 $\bm a$ 满足一个非线性协商方程 $\bm H\bm a=\bm a^{-1}$。这说明 Nash 方向里各任务的权重是被“协商关系”共同决定的，不是谁随手指定的加权。数值求解有条件：当各任务梯度线性无关时，Gram 矩阵 $\bm H$ 正定，方程可稳定求解；但如果梯度秩亏、近乎线性相关、或压根不存在严格共同下降方向，$\bm H$ 会奇异或近奇异，这时不能不加检查就去求逆。

原论文使用凸代理与凸凹过程近似求权重。这里为手算另给一个教学求法，不声称复现其实现。在$\bm H\succ0$、$\bm a>0$下，考虑二次项与对数屏障的和，其梯度为零正好是式~\eqref{eq:ket-sharing-nash-system}：
\begin{equation}
  \psi(\bm a)
  =\tfrac12\bm a^{\top}\bm H\bm a
   -\sum_k\log a_k,
  \qquad a_k>0.
  \label{eq:ket-sharing-nash-newton}
\end{equation}
该函数严格凸，梯度$\bm r=\bm H\bm a-\bm a^{-1}$，海森矩阵$\bm H+\operatorname{diag}(a_k^{-2})$正定。以全正向量初始化，求解牛顿方向的线性方程：
\[
  \bigl[\bm H+\operatorname{diag}(a_k^{-2})\bigr]
  \bm\delta=-\bm r
\]
所得$\bm\delta$是系数增量。通过回溯选择步长，使全部$a_k$继续为正且$\psi$减小；以方程残差达到预设容差为停止条件，再归一化$\bm G\bm a$。求解误差存在时，这个方向也只是协商解的近似。

重估系数需要取得逐任务梯度、形成Gram矩阵并求解$K$维子问题。复用多步前的系数可省重估费用，但不再精确满足当前梯度的协商方程。

沿用零点教学梯度，$\bm H=\left[\begin{smallmatrix}1&-2\\-2&8\end{smallmatrix}\right]$。两行方程分别乘$a_1,a_2$，得$a_1(a_1-2a_2)=1$和$a_2(-2a_1+8a_2)=1$；相减后$a_1^2=8a_2^2$。选择全正根便有
\begin{align}
  a_2&=(8-4\sqrt2)^{-1/2},
  \qquad a_1=2\sqrt2\,a_2,
  \nonumber\\
  \bm d_{\mathrm N}&\approx(0.3827,0.9239)^{\top}.
\end{align}
前面MGDA方向长度不是1，不能直接与单位Nash方向比较收益。将其归一化为$\bm d_{\mathrm M}=(2,3)^{\top}/\sqrt{13}$，再对同一梯度计算两种单位方向的效用：
\begin{align}
  (u_1,u_2)(\bm d_{\mathrm M})
    &\approx(0.5547,0.5547),
  \nonumber\\
  (u_1,u_2)(\bm d_{\mathrm N})
    &\approx(0.3827,1.0824).
  \label{eq:ket-sharing-nash-numbers}
\end{align}
算出来的对比很能说明“不同准则＝不同的公平观”：把 MGDA 方向归一化到单位长度后，它给两个任务完全相等的效用（约 0.5547、0.5547），像“平均主义”；Nash 方向则给任务 1 较小的正效用（约 0.3827）、给任务 2 较大效用（约 1.0824）。Nash 的收益乘积（约 0.4142）确实大于归一化 MGDA 的乘积（约 0.3077），所以从“乘积最大化/协商”标准看它更优；但它并不是在每一项上都超过 MGDA（任务 1 反而拿得更少）。结论是：这些优化准则选择的是不同的取舍点，并不存在一个在每一维收益上都一致占优的排名，用哪个取决于你认同哪种分配哲学。

图~\ref{fig:ket-sharing-gradient-geometry}固定同一坐标与原梯度，分别画投影后的组合、凸包最近点和协商单位方向。观察箭头时应区分梯度组合$\bm d$与模型参数移动：优化沿$-\eta\bm d$执行，不能把图中箭头本身当作实际位移。

\begin{figure*}[htbp]
  \centering
  % Teaching geometry from eq:ket-sharing-toy-losses; no experimental data.
% Source gradients are g1=(1,0), g2=(-2,2), not (-2,3).
% All panels use 11.5 mm per coordinate unit; canvas is 169 x 74 mm.
% Arrows show gradients or their combinations d; displacement is -eta*d.
\begin{tikzpicture}[
  x=1mm,y=1mm,
  font=\figurefont\footnotesize,
  text=BookInk,
  >={Latex[length=1.7mm,width=1.2mm]},
  axis/.style={draw=BookMuted,line width=0.6pt,-{Latex}},
  tick/.style={draw=BookMuted,line width=0.5pt},
  original/.style={draw=FigureInput,line width=0.9pt,-{Latex}},
  projected/.style={draw=FigureModel,line width=1pt,-{Latex}},
  result/.style={draw=FigureLoss,line width=1.1pt,-{Latex}},
  guide/.style={draw=BookMuted,dashed,line width=0.7pt},
  leader/.style={draw=BookMuted,line width=0.5pt},
  heading/.style={anchor=west,font=\figurefont\footnotesize\bfseries},
  label/.style={inner sep=0.6mm},
  note/.style={align=center,inner sep=0pt}
]
  \path[use as bounding box] (0,0) rectangle (169,74);
  \node[heading] at (0,70) {(a) PCGrad：投影};
  \node[heading] at (58,70) {(b) MGDA：凸包最短点};
  \node[heading] at (116,70) {(c) 单位方向与共同下降};

  % Identical origins, ranges and physical scales in all three panels.
  \foreach \offset in {0,58,116} {
    \begin{scope}[shift={(\offset,0)}]
      \begin{scope}[shift={(29,25)},x=11.5mm,y=11.5mm]
        \draw[axis] (-2.25,0) -- (1.45,0);
        \draw[axis] (0,-1.2) -- (0,2.9);
        \foreach \v in {-2,-1,1} {
          \draw[tick] (\v,-0.045) -- (\v,0.045);
          \node[label,anchor=north] at (\v,-0.09) {$\v$};
        }
        \foreach \v in {-1,1,2} {
          \draw[tick] (-0.045,\v) -- (0.045,\v);
          \node[label,anchor=east] at (-0.09,\v) {$\v$};
        }
        \node[label,anchor=north east] at (-0.06,-0.06) {$0$};
      \end{scope}
      \node[label,anchor=west] at (46.5,25) {$d_1$};
      \node[label,anchor=south] at (29,60) {$d_2$};
    \end{scope}
  }

  \begin{scope}[shift={(29,25)},x=11.5mm,y=11.5mm]
    \coordinate (p1) at (0.5,0.5);
    \coordinate (p2) at (0,2);
    \draw[guide] (1,0) -- (p1);
    \draw[guide] (-2,2) -- (p2);
    \draw[original] (0,0) -- (1,0);
    \draw[original] (0,0) -- (-2,2);
    \draw[projected] (0,0) -- (p1);
    \draw[projected] (0,0) -- (p2);
    \draw[result] (0,0) -- (0.5,2.5);
    % The two removed components are orthogonal to the projected vectors.
    \draw[leader] (0.42,0.42) -- (0.50,0.34) -- (0.58,0.42);
    \draw[leader] (-0.14,2) -- (-0.14,1.86) -- (0,1.86);
  \end{scope}
  \node[label,anchor=west] at (42,29) {$\bm g_1$};
  \node[label,anchor=south] at (6,49) {$\bm g_2$};
  \node[label,anchor=west] at (38,35) {$\widetilde{\bm g}_1$};
  \draw[leader] (p1) -- (38,34);
  \node[label,anchor=south] at (21,54) {$\widetilde{\bm g}_2$};
  \draw[leader] (p2) -- (21,54);
  \node[label,anchor=west] at (37,56) {$\bm d_{\mathrm{PC}}$};
  \node[note] at (26.5,4) {$\bm d_{\mathrm{PC}}=(0.5,2.5)^{\top}$};

  \begin{scope}[shift={(87,25)},x=11.5mm,y=11.5mm]
    \coordinate (mgda) at ({4/13},{6/13});
    \draw[draw=BookInk,line width=1pt] (-2,2) -- (1,0);
    \fill[FigureInput] (-2,2) circle[radius=0.9mm];
    \fill[FigureInput] (1,0) circle[radius=0.9mm];
    \draw[projected] (0,0) -- (mgda);
    \fill[FigureModel] (mgda) circle[radius=0.7mm];
    % Perpendicularity to g1-g2=(3,-2) identifies the closest point.
    \draw[leader]
      ({4/13+0.16*3/sqrt(13)},{6/13-0.16*2/sqrt(13)})
      -- ({4/13+0.16/sqrt(13)},{6/13-0.16*5/sqrt(13)})
      -- ({4/13-0.16*2/sqrt(13)},{6/13-0.16*3/sqrt(13)});
  \end{scope}
  \node[label,anchor=south] at (64,49) {$\bm g_2$};
  \node[label,anchor=west] at (99,29) {$\bm g_1$};
  \node[label,anchor=west] at (96,37) {$\bm d^\star$};
  \draw[leader] (mgda) -- (96,36);
  \node[note] at (84.5,4) {$\bm d^\star=(4/13,6/13)^{\top}$};

  \begin{scope}[shift={(145,25)},x=11.5mm,y=11.5mm]
    % g1.d>0 and g2.d>0 give d1>0 and d2>d1: angles 45 to 90 deg.
    \fill[FigureModel!12] (0,0) -- (45:1)
      arc[start angle=45,end angle=90,radius=11.5mm] -- cycle;
    \draw[draw=BookMuted,dotted,line width=0.7pt] (0,0) circle[radius=11.5mm];
    \draw[guide] (45:1) -- (0,0) -- (90:1);
    \coordinate (unitm) at ({2/sqrt(13)},{3/sqrt(13)});
    \coordinate (unitn) at ({sin(22.5)},{cos(22.5)});
    \draw[projected] (0,0) -- (unitm);
    \draw[result,dashed] (0,0) -- (unitn);
  \end{scope}
  \node[note,anchor=north west,align=left] at (118,54)
    {共同下降区域\\$d_1>0$\\$d_2>d_1$};
  \node[label,anchor=south] at (151,44) {$\bm d_{\mathrm N}$};
  \draw[leader] (unitn) -- (151,44);
  \node[label,anchor=west] at (159,38) {$\bm d_{\mathrm M}$};
  \draw[leader] (unitm) -- (159,38);
  \node[label,anchor=east] at (133,12) {单位圆};
  \draw[leader] (134,13) -- (138.1,15.8);
  \node[note] at (142.5,4) {实际位移：$-\eta\bm d$};
\end{tikzpicture}

  \caption{固定原梯度的二维教学对照，各面板采用相同坐标尺度，$\bm g_1=(1,0)^{\top}$、$\bm g_2=(-2,2)^{\top}$。(a)虚线从原梯度端点连到PCGrad投影端点，两个投影向量求和为$\bm d_{\mathrm{PC}}$。(b)原梯度端点之间的线段是凸包，距原点最近的是$\bm d^\star=(4/13,6/13)^{\top}$。(c)点线圆限制单位方向，浅色扇区示意$d_1>0,d_2>d_1$的共同一阶下降组合，径向虚线边界不属于严格下降区域。实线和虚线分别是单位MGDA方向$\bm d_{\mathrm M}$与Nash方向$\bm d_{\mathrm N}$。各箭头均为梯度组合，模型参数位移需反向取$-\eta\bm d$；图不表示最终任务性能次序。}
  \label{fig:ket-sharing-gradient-geometry}
\end{figure*}
\FloatBarrier

Nash 协商解还附带一个著名性质——比例公平（proportional fairness）：在零参考效用、同一可行集下，对任何其他可行的效用组合 $u_k$，按最优效用 $u_k^\star$ 归一化后的相对变化之和 $\sum_k(u_k-u_k^\star)/u_k^\star$ 都不为正。换句话说，你无法在不引起别人相对损失更大的情况下，让某些人相对获益。它还有个尺度不变性：把某项损失乘一个正常数，只给对数目标加个常量，不改变最优的单位方向。但要警惕过度解读：这只是局部效用上的数学尺度性质，它既不能替代人为设定的任务优先级，也绝不意味着社会学意义上的“人群服务公平”。各群体的成本与收益必须另设评价单位或约束来考察，不能用“总损失下降”去冒充“每个任务、每个群体都受益”。

面对这么多精巧的协调器，必须让它们接受“简单基线”的实证检验，否则容易高估其价值。Xin 等人在机器翻译上做了一组有警示意义的对照：固定英语到法语为一个任务，分别搭配英语到汉语、德语或罗马尼亚语，比较 PCGrad、MGDA、GradNorm、他们提出的 IMTL-G，以及最简单的“随机损失加权”。评价用 BLEU（衡量译文和参考译文的片段匹配，越高越好）。结果发现：在他们所用的 Transformer、Adam 和调参范围内，这些专用协调器最终给出的分数取舍，并没有超出“简单地扫描一下任务采样比例”就能达到的取舍；而且一旦改变评价偏好，方法之间的排名还会跟着改变\scite{ket-xin2022-mto}。当然，负面证据有其适用边界（该研究没测 Nash-MTL，任务也限于翻译），不能据此否定所有协调方法，但它提醒我们：在上复杂协调器之前，先把采样比例这种便宜基线调好。

前面所有方法最终都训出一组确定的参数，对应帕累托前沿上的一个固定点。可如果不同植物园对“物种错误”和“病害漏检”的取舍各不相同、而且要到上线时才告诉你，为每一种偏好都从头训一个模型显然不现实。\term{帕累托超网络}（Pareto HyperNetwork, PHN）用一个巧妙的办法一次训练、按需生成：它用一个超网络 $H_{\bm\psi}$ 把偏好向量 $\bm r$（$r_k\geq0$、和为 1，描述“多在乎每个任务”）直接映射成一整套任务模型参数 $\bm\theta=H_{\bm\psi}(\bm r)$，$\bm\psi$ 是这个生成网络的参数。训练时随机抽取各种偏好，用对应的加权损失或偏好协调方向去更新 $\bm\psi$，逼它学会整条前沿；部署时植物园只需声明偏好，超网络即时生成相应模型。这里联合训练的是“偏好到模型参数”这个映射；它和第~\ref{sec:ket-integration-deltas}节从已成形的参数增量构造偏好相关融合是不同思路。它的局限也要讲清：训练时偏好的覆盖范围、超网络自身的容量、优化误差都会限制能触达的取舍，绝不能把超网络训到的某个驻点当成“完整真实风险前沿”的全局保证；评价时仍要报告若干预定偏好下的逐任务结果，以及生成模型的额外开销\scite{ket-navon2021-phn}

最后划清两个边界。其一，上面所有协调处理的都是“参与了训练的任务”；要检验一个从未共同训练过的新任务是否也能受益，必须像第~\ref{sec:ket-transfer-task-splits}节那样把整个任务留出（而不只是留出某些已见任务的图像），并规定它能用哪些适应证据。其二，如果任务冲突长期、反复地持续，不要只在更新方向层面反复打补丁，而应回到更上游去检查分组是否合理、共享深度是否过深、输出关系规则是否设错——优化器不是万能的：如果当前结构压根没有为某个特征预留容量，那么再聪明的梯度协调也无法在这个结构内部把丢失的特征变回来。

\section{联合训练资源的分配}
\label{sec:ket-sharing-resources}

到这里，目标函数看起来已经定义好了，但“写下同一个联合目标”并不等于“各方案拿到的证据、花掉的资源相同”。资源分配是另一层常被忽略、却足以决定成败的设计：花钱补标一批病害，增加的是监督信息量；提高病害任务被抽中的概率，增加的是它使用已有样本的次数；把专家网络扩大，增加的则是算力和表示自由度——这三件事都能让病害表现变好，但它们是三种完全不同的投入，不能用一句“都加强了病害”糊弄过去。本节就分别核算三类资源：样本与标签、训练更新机会、计算与存储预算，让实际训练真正对应得上前面定下的要求。只比较“总训练步数”是不够的，因为同样 1000 步，证据构成和花费可能天差地别。

\subsection{样本与标签资源的分配}
\label{sec:ket-sharing-labels}

手里只有有限的标注预算时，先想清楚“标在哪”比“怎么训”更重要。有两种截然不同的花法：一是去覆盖更多植株（增加独立对象的多样性，让模型见到更广的世界）；二是给已经拍过的图像补齐其他任务的标签（让同一植株的类别、病害、位置彼此对应，便于学习任务间的输出关系）。三种常见策略——按任务均分、优先补最薄弱的任务、按对象成套标注（一张图同时标三类）——会改变标签的分布和可用的跨任务对应关系。所以“标签总数相同”根本不能保证提供相同的监督依据：同样 1000 个标签，覆盖 1000 株各标一个任务，和覆盖 333 株每株标三个任务，学到的东西很不一样。

现实里并不是每张图都有三类标签，于是“标签缺失时怎么算损失”成为第一个具体问题。最重要的原则是：只对真实观察到的标签计算监督，绝不能给缺失标签偷偷填默认值。用掩码 $m_{i,k}\in\{0,1\}$ 标记样本 $i$ 是否有任务 $k$ 的可用标签，$n_{B,k}=\sum_{i\in B}m_{i,k}$ 是当前批次 $B$ 中该任务实际可见的标签数。只在 $n_{B,k}>0$ 时，对这些可见标签求平均：
\begin{equation}
  L_k(B)=
  \frac{1}{n_{B,k}}
  \sum_{i\in B}m_{i,k}
  \ell_k(f_k(\bm x_i),y_{i,k}).
  \label{eq:ket-sharing-label-mask}
\end{equation}
代码里常见的“让缺失项乘零”只是一种掩码写法，含义必须守住：程序实际上只能读取 $m_{i,k}=1$ 的那些标签；缺失时的 $y_{i,k}$ 是个占位符，绝不能被当成有效的输入或答案。若 $n_{B,k}=0$（这一批里该任务一个标签都没有），分母为零、平均无定义，本批就该跳过该任务、不产生监督，而不是硬算。还有两个语义错误必须杜绝：病害状态“未知”不等于“健康”，类别“未知”也不等于“负类”——把缺失当否定，会系统性教给模型错误的规律。

一个数字例子能说清“分母选择”如何悄悄改变任务权重。设批次里图像 1 只有类别标签，图像 2 三类齐全，图像 3 只有病害和定位；各任务两项可见损失分别为：类别 $(0.2,0.6)$、病害 $(0.8,0.4)$、定位 $(0.3,0.5)$。按可见数平均（正确做法），三项损失都是 $0.4$、$0.6$、$0.4$，等权对待。但若不分青红皂白都除以完整批次大小 3，就等于给每项乘上了各自的可见率 $2/3$；更一般地，当各任务可见数不同时，这个乘数也不同——标签少的任务被自动缩小了权重，且缩小幅度隐蔽、随批次变化。两种平均都可以作为一种事先声明的目标约定，但绝不能在不说明的情况下无声互换，否则任务的相对作用就失控了。

只对可见标签求平均，技术上没问题，但它带来一个更隐蔽的统计问题：删除缺失项并不会消除选择偏差。设想专家总是优先标注那些病斑明显的图像（健康的、模糊的就懒得标），那么可见标签其实来自一个被“是否值得标注”筛选过的条件总体 $m_k=1$，你算的可见平均反映的是 $\mathbb E[\ell_k\mid m_k=1]$（有标签的样本上的损失），而不是一般意义上、面向所有图像的总体风险。要恢复式~\eqref{eq:ket-sharing-task-risk}真正想要的总体对象，需要对“标注是怎么产生的”做出假设并加以校正。先固定任务和预测器，假设批次从目标输入总体独立同分布抽取，再有一个标注机制决定揭示哪些答案。

记固定预测器下样本 $i$ 在任务 $k$ 上的损失为 $\ell_{i,k}$。引入两个关键假设：其一，给定输入后，“是否被标注”与未知答案条件独立 $m_{i,k}\perp y_{i,k}\mid\bm x_i$（即标注与否只取决于看得见的图像特征，不取决于那个还没人知道的答案）；其二，可见概率（倾向性）$e_k(\bm x)=\mathbb P(m_{i,k}=1\mid\bm x_i=\bm x)$ 已知，且在目标分布的支撑上处处为正（任何图像都有非零机会被标注）。这时就可以用逆概率加权（IPW）补偿那些“很少被标注”的输入——越不容易被标的，一旦被标到就给越大权重。在损失可积、批次非空时，以完整批次大小为分母构造无偏估计：
\begin{equation}
  \widehat R_k(B)=\frac{1}{|B|}
    \sum_{i\in B}\frac{m_{i,k}}{e_k(\bm x_i)}\ell_{i,k}.
  \label{eq:ket-sharing-label-ipw}
\end{equation}
为什么这样加权就能还原总体风险？关键是验证它的无偏性。先固定输入和答案不动，只对“标注是否出现”这个随机事件取条件期望，看看加权项的期望是否恰好等于原始损失：
\begin{align}
  \mathbb E\!\left[
    \frac{m_{i,k}\ell_{i,k}}{e_k(\bm x_i)}
    \,\middle|\,\bm x_i,y_{i,k}\right]
  &=\frac{\ell_{i,k}}{e_k(\bm x_i)}
    \mathbb P(m_{i,k}=1\mid\bm x_i,y_{i,k})
  \nonumber\\
  &=\ell_{i,k}.
  \label{eq:ket-sharing-label-ipw-expectation}
\end{align}
第二个等号正是用了“给定输入后标注与答案独立”的假设（$\mathbb P(m=1\mid\bm x,y)=\mathbb P(m=1\mid\bm x)=e_k(\bm x)$，于是 $e_k$ 恰好约掉）；之后再对输入和答案取期望，就得到 $\mathbb E[\widehat R_k(B)]=R_k(f_k)$——这个加权平均在期望意义下无偏。要理解实现：没被标注的损失在这里只是一个理论上的随机变量（乘了 $m=0$ 就是零），程序实际仍只累加可见项，并不需要凭空造标签。但结论有严格前提：若 $e_k$ 是估计出来而非已知、若可见性其实依赖未观察到的答案（独立性破裂）、或某类图像的标注概率为零（支撑不覆盖），无偏所需的依据就缺失了，不能照搬这个漂亮结论。

做了逆概率加权后，分母用哪个也必须与目标总体对应，否则又会偏。用一个例子看：设批次有四项损失都为 1，每项可见概率都是 $1/2$，恰好观察到两个标签。加权求和为 $2/(1/2)=4$；若除以完整批次大小 4，正确得到 1（还原总体均值）；可若除以可见数 2，却得到 2（高估了）。实践中还有一种更稳的变体——\term{自归一化估计}（self-normalized estimator），它不除以固定的批次大小，而是除以本批实际出现的权重之和，让估计自动按实际权重总量定尺度：
\begin{equation}
  \widehat R_{k,\mathrm{SN}}(B)
  =\frac{\sum_{i\in B}m_{i,k}\ell_{i,k}/e_k(\bm x_i)}
    {\sum_{i\in B}m_{i,k}/e_k(\bm x_i)}.
  \label{eq:ket-sharing-label-self-normalized}
\end{equation}
只有当权重和（分母）为正时这个比值才能计算。代价是：分母本身是随机的（依赖这批恰好抽到哪些可见标签），所以自归一化估计在有限样本下一般是有偏的，只是方差往往更小、更稳定——这是偏差和方差之间的交换。上面那个常数损失例子里它恰好也得到 1，但那只是“损失全相等”的特殊情形，不能据此认为它和完整批次的无偏估计在一般情况下恒等。

“一个标签都没有的批次”是检验实现是否正确的绝佳试金石，因为三种量在此时的行为完全不同：完整批次的 IPW 估计取值为 0——但这个 0 表示“本次随机梯度贡献为空”，绝不表示“真实损失就是 0”；而可见平均和自归一化比值在此时分母为零、本应是未定义。工程上本批可以跳过该任务的监督反传，但风险平均仍要保留完整批次的计数，预先约定的任务权重也不能因为“这批某任务碰巧没标签”就重新归一化（否则目标被随机批次改变）。一个理论细节：固定模型时，若简单地把零贡献批次整批删掉、只在非空批次上平均，得到的其实是 $R_k(f_k)/\mathbb P(n_{B,k}>0)$ 这个条件期望，已经偷偷改变了风险口径。

还要把“用于调权的学习进度”和上面这个“随机损失贡献”严格分开：前者是跨阶段的、相对稳定的量，后者是单批次随标签缺失随机为零的量，绝不能拿后者当前者用。GradNorm 这类方法在缺标签时应保留最近一次有效的进度值，而不是读到零；一种保守而稳妥的处理是：本批干脆暂缓整组的权重调整、沿用原有权重，但各任务的监督梯度仍按选定的（比如 IPW）估计式照常计算。异步更新、或只对本批可见任务调权，都需要另行制定规则，底线是不能把“这个随机的零贡献”误记成“这个任务已经学好了、损失为零”。最后还有个严谨性提醒：风险估计无偏，只有在期望和求导可交换时才给出无偏梯度；它更不能直接外推到依赖随机状态的动态调权更新里（那条链路上无偏性未必保持）。相关的分布偏移校正思路见第~\ref{sec:ket-update-weighting}节。

核算样本费用时，要按“实际证据”计数，而不是按文件数。给同一张图标三类标签，得到的是三项监督，但它并没有因此变成三张相互独立的照片；同一张原图的多个增强版本、同一植株的连续拍摄，也都必须和“独立对象”区分开——它们携带的信息高度重叠。按植株划分训练/测试时，同一植株的所有样本要整组隔离，不能一个增强进训练、一个进测试。还要注意，标签的查询与核验属于第 5 章“监督信号获取”的成本，本节只管“已有标签如何使用”；在训练中把同一条样本重复抽取多次，并不等于新增了证据，只是重复利用了已有的那一份信息。

将任务改成自然语言接口时，也要保留这些单位。T0用多个模板描述样本集对应的任务，并以留出任务类别检查未参与共同训练的要求\scite{ket-sanh2022-t0}。未见问法、未见样本来源和未见预测类别分别改变措辞、样本集和任务范围；一种留出不能代替另外两种。报告任务泛化前，应明确留出了哪一层。

FLAN的扩展研究混合多种指令训练样本，并讨论任务、模板与来源重复的核算\scite{ket-chung2024-flan}。同样，把“是否患病”改写多次不应计成多项独立能力。合并时分别检查对象、改写预测输入和答案，防止重复内容跨入评价样本集，并说明增加的是任务要求、独立观测还是表达形式。

\subsection{训练更新机会的分配}
\label{sec:ket-sharing-sampling}

标签清单完全不变，光改变“每个任务被抽到多少次”，也能显著改变训练结果——这就是更新机会（采样）的分配。为了把“采样概率”和“损失权重”如何共同起作⽤算清楚，先考察最简单的情形：权重固定、线性加权更新。要提醒的是，PCGrad、MGDA、Nash-MTL 这类方法要求在同一模型状态下同时取得所有任务的梯度；如果你改成各任务分步收集、或复用陈旧梯度，那是另外的近似，下面基于线性期望得到的干净结论不能直接套用。

设定每一步先以概率 $p_k>0$ 抽取任务 $k$，再从该任务的已标注样本里均匀抽一个小批次，最后用权重 $w_k$ 给抽到的梯度加权。固定当前参数、概率和权重，对“抽哪个任务”和“任务内抽哪些样本”两层随机性都取期望，单步梯度的期望为：
\begin{equation}
  \mathbb E[\widehat{\bm g}]
  =\sum_{k=1}^{K}p_kw_k
    \nabla\hat R_k.
  \label{eq:ket-sharing-effective-objective}
\end{equation}
这个等式揭示了一个关键事实：优化器实际看到的期望目标是 $\sum_k p_kw_k\hat R_k$，真正起作用的系数是“采样概率 × 损失权重”的乘积，而不只是你在损失代码里写的那个 $w_k$。换句话说，就算你给所有任务设了相等的损失权重，只要采样概率不均等，实际优化的目标依然是倾斜的。推论是：如果任务内部还按难度做了偏置采样，那条件样本分布也跟着变了，必须继续校正或相应改写目标，才能得到你想要的期望。

两种自然的采样选择对应两种不同的“公平观”，要分清。令 $n_k$ 为任务 $k$ 的可抽样本数：取 $p_k=n_k/\sum_j n_j$（按样本量比例采样）、再配等权损失，等价于直接对全部样本的样本平均优化——数据多的任务自然主导；取 $p_k=1/K$（各任务等概率采样）则等价于“任务平均”——不管数据多少，每个任务话语权相同。这是两种不同的目标，没有对错，取决于你想对样本负责还是对任务负责。$T$ 步里任务 $k$ 被抽中次数的期望是 $Tp_k$（实际有随机波动）。还要点破一个常见误解：提高采样频率是让任务“更多次产生梯度”，提高损失权重只是“把已经抽到的梯度放大”，两者机制不同；而且它们都不会凭空增加独立标签的数量——小样本任务被重复抽，看的还是那点数据。

这里有个灵活的工程需求：你心里想要的目标系数是 $q_k$，但出于某种原因（比如设备效率、负载均衡）实际采样频率只能是 $p_k$，二者不一致怎么办？用重要性采样就能解耦——先以 $q_k\geq0,\sum_kq_k=1$ 指定“目标里各任务应占的系数”（注意这里的 $q_k$ 和 GradNorm 里的损失进度比只是记号相同、含义无关）；当抽到任务 $k$ 时，给它的梯度乘上补偿因子 $q_k/p_k$，得到：
\begin{equation}
  \widehat{\bm g}
  =\frac{q_k}{p_k}\nabla\ell_k,
  \qquad
  \mathbb E[\widehat{\bm g}]
  =\sum_kq_k\nabla\hat R_k.
  \label{eq:ket-sharing-importance-sampling}
\end{equation}
无偏的条件是：每个 $q_k>0$ 的任务都必须有正的抽取概率 $p_k>0$（否则永远补偿不到，权重趋于无穷），并且在这一步估计里把 $q_k/p_k$ 当作常量、不对采样决策本身求导。采样概率可以根据训练历史动态调整，只要“给定历史后”上述条件仍满足，相应的条件期望结论就依然成立。如果同时还设了损失权重，要核对最终生效系数是“概率补偿 × 损失权重”的连乘积，避免对同一差异补偿两次。

重要性采样能让不同采样方案有相同的期望目标，但“期望相同”不等于“训练一样稳”——它们的方差可能差很多。设任务 $k$ 内随机样本梯度为 $\bm g_k$、均值为 $\bm\mu_k=\mathbb E[\bm g_k]$、二阶矩有限。把“校正后梯度相对其均值的偏差”的平方范数展开，可得到方差表达式：
\begin{align}
  &\mathbb E\left\|
    \widehat{\bm g}-\sum_kq_k\bm\mu_k
  \right\|_2^2
  \nonumber\\
  &\quad=
  \sum_k\frac{q_k^2}{p_k}\mathbb E\|\bm g_k\|_2^2
  -\left\|\sum_kq_k\bm\mu_k\right\|_2^2.
  \label{eq:ket-sharing-sampling-variance}
\end{align}
式中二阶矩那一项带因子 $1/p_k$：当某任务采样概率极小，它偶尔被抽中时补偿因子 $q_k/p_k$ 会非常大，导致梯度剧烈跳动、方差激增。所以采样概率不是越小越省、也不是盲目给薄弱任务增频就好，而要综合权衡所有任务的二阶矩和单步成本；而且别忘了，给一个小样本池增频，本质仍是在重复同样的证据、并不增加新信息。控制方差有两类常用手段：给采样概率设一个正下界（限制补偿倍数），或直接截断过大的权重——但截断通常会引入偏差，属于“用偏差换稳定”，必须如实记录这一改变。

除了固定的静态分配，还可以让采样随训练进程变化：一种是预先排定课程（curriculum），比如早期多练较容易、数据多的类别任务，后期再逐步增加病害和定位的比重；另一种是根据验证集上“哪些任务还没达标”的反馈动态加抽。评价这类方法时要固定总更新次数或总算力预算、记录各阶段的实际采样次数和校正方式，并且最终测试集绝不能参与这种反馈选择。一个概念上的提醒：未经重要性校正的动态采样，其实在训练过程中不断改变着有效目标本身，所以不能把它轻描淡写地说成“只是在更高效地求解同一个固定目标”——目标是随时间漂移的。

\subsection{计算与存储预算的分配}
\label{sec:ket-sharing-budget}

先澄清计算和存储是两本不同的账，别混为一谈。计算方面：当同一图像在一次请求里同时要类别、病害、位置时，共享骨干确实只需前向一次，三个头复用同一份特征；但不同图像之间不能复用——来一张新图就得重算一次骨干。存储方面则按参数算：设骨干参数量为 $P_{\mathrm b}$、每个头为 $P_{\mathrm h}$，三个完全独立的模型要存 $3(P_{\mathrm b}+P_{\mathrm h})$，而“一个共享骨干 + 三个头”只需 $P_{\mathrm b}+3P_{\mathrm h}$。共享省下的是“重复骨干”那部分参数和与之对应的部分重复计算，但这绝不等于所有调用场景都等比例省钱：如果上线后其实只服务类别一个任务，那为病害、定位多训的头和多花的训练费用，就得另外说明它们带来的间接价值（比如帮助骨干学到了更好的特征）。

用一组教学数字（单位百万）把账算具体：取 $P_{\mathrm b}=12$、$P_{\mathrm h}=1$。三个模型分别训练要存 $3\times(12+1)=39$；共享骨干加三个头只存 $12+3=15$；而真正跑某个单任务时，前向路径只涉及骨干加它自己的头，即 13。这个对比只算存储、不含预测质量的高下。这里还有个常被忽略的公平性问题：如果你给共享方案也放宽到 39 的总参数预算，省下的重复骨干容量可以反过来把骨干扩大到 36（加三个头仍是 39），或用余钱增设私有分支——这种“等参数预算下把骨干加宽”的收益必须实测，不能默认。反过来也要记住：参数量相同并不保证浮点运算量、推理时延、训练显存都相同（一个宽而浅、一个窄而深可以同参数但速度差很多）。

条件专家（如 MMoE）这类结构让“存多少”和“算多少”进一步分离，需要两个指标分别记账。\term{总容量}（total capacity）指磁盘/显存里实际保存的全部参数（所有专家都得存着）；\term{激活容量}（active capacity）指一次请求中真正参与计算的参数（只算被门控选中执行的专家）。设共享前端为 $P_0$、每个专家 $P_{\mathrm e}$、任务 $k$ 的门控和头分别为 $P_{\mathrm g,k},P_{\mathrm h,k}$，并用 $A_k(\bm x)$ 表示该输入下真正被执行的专家集合，两种容量写为：
\begin{align}
  P_{\mathrm{total}}
    &=P_0+EP_{\mathrm e}
      +\sum_k(P_{\mathrm g,k}+P_{\mathrm h,k}),
  \nonumber\\
  P_{\mathrm{active},k}
    &=P_0+|A_k(\bm x)|P_{\mathrm e}
      +P_{\mathrm g,k}+P_{\mathrm h,k},
  \label{eq:ket-sharing-expert-budget}
\end{align}
标准（稠密）MMoE 通常把 $E$ 个专家全部执行，即 $|A_k|=E$，此时激活容量并不比总容量小；只有额外引入稀疏选择机制、能真正跳过未选中专家，$|A_k|$ 才会降下来。多任务同读一张图时还有缓存效应：若专家结果可复用，实际需要执行的专家是各任务所选集合的并集 $\bigcup_k A_k(\bm x)$，而不是把每个任务的开销简单相加（否则重复计了被多任务共用的专家）。最要紧的一个区分是：没被执行的专家参数依然躺在设备存储里占地方，所以“激活参数少、单次算得快”绝不等于“模型驻留占用少”——这正是要分设两个容量指标的原因。

数字演练：取 $P_0=4,E=4,P_{\mathrm e}=6$，三个门控各 0.1、头各 1（百万）。总容量为 $4+4\times6+3\times(0.1+1)=31.3$。单个任务走稠密路径（执行全部 4 个专家）时激活 $4+4\times6+0.1+1=29.1$；若加上稀疏机制、每次只执行一个专家，该路径降到 $4+6+0.1+1=11.1$，但磁盘上仍要保存全部 31.3。这组数字把“算得少”和“存得少”的差别摆得很清楚。另外，参数计数离真实开销还差很多：门控网络的计算、多个专家输出的聚合、批次形状（batch）、以及跨设备的参数传输都要花时间和带宽，所以参数量不能直接换算成推理时延。

训练时的显存账还要再加上梯度和优化器状态，往往比模型本身大好几倍。若 $P$ 个可训练参数都用 32 位浮点（FP32，每数 4 字节），Adam 优化器要分别保存：模型参数、梯度、一阶动量、二阶动量，共四份数组，约 $16P$ 字节（这还没算前向激活和各种缓冲）。举例：15 百万可训练参数，光这四份就约 $240\,$MB（其中单份模型参数 $60\,$MB；此处 MB 按 $10^6$ 字节计）。稀疏训练也要小心：某些专家这一步不活跃、不产生梯度，但优化器通常仍为它们保留历史动量状态，这部分存储省不掉。混合精度训练、激活重计算、优化器状态分片（sharding）、参数卸载（offload）等技术会改变这笔账；而批次大小、输入分辨率和中间激活才常常是显存峰值的真正大头，所以最终必须按实际执行配置去实测，而不能只拿参数量估算。

表~\ref{tab:ket-sharing-resource-controls}固定类别、病害和定位三项教学样本集，分别有$600,300,100$条样本标注；每步抽一任务的一项样本标注，共1000步，训练样本、优化器与损失单位不变。$C_0$表示基准共享容量，容量变化不预设计算线性变化。读表可见，改权重与改采样能得到相同的任务平均目标，却保留不同机会和方差；改容量则改变候选模型及单步费用。

\begin{table*}[htbp]
  \centering
  \small
  \begin{tabularx}{\linewidth}{p{18mm}p{36mm}p{25mm}p{28mm}X}
    \toprule
    教学方案 & 抽取概率$\bm p$及权重$\bm w$ & 期望目标系数$p_kw_k$ & 抽取次数的期望 & 模型容量与费用核算 \\
    \midrule
    基准 &
    $\bm p=(0.6,0.3,0.1)$\newline
    $\bm w=(1,1,1)$ &
    $(0.6,0.3,0.1)$ &
    $(600,300,100)$ &
    容量为$C_0$；优化样本标注平均，单步费用作基准。 \\
    只改权重 &
    $\bm p$同基准\newline
    $\bm w=(5/9,10/9,10/3)$ &
    $(1/3,1/3,1/3)$ &
    $(600,300,100)$ &
    容量仍为$C_0$；抽取次数不变，低频任务的梯度放大。 \\
    只改采样 &
    $\bm p=(1/3,1/3,1/3)$\newline
    $\bm w=(1,1,1)$ &
    $(1/3,1/3,1/3)$ &
    各任务$1000/3$ &
    容量仍为$C_0$；小样本任务重复增多，按各任务单步费用计总计算。 \\
    只改容量 &
    $\bm p,\bm w$同基准 &
    $(0.6,0.3,0.1)$ &
    $(600,300,100)$ &
    容量增为$2C_0$；抽取与目标不变，重新测量单步成本。 \\
    联合改变 &
    $\bm p=(0.4,0.3,0.3)$\newline
    $\bm w=(5/6,10/9,10/9)$ &
    $(1/3,1/3,1/3)$ &
    $(400,300,300)$ &
    容量为$2C_0$；权重补偿抽取概率，分别核算机会与扩容。 \\
    \bottomrule
  \end{tabularx}
  \caption{相同任务训练样本下改变权重、抽取概率与容量的教学对照，数字不含实测性能，向量顺序为类别、病害、定位。抽取次数表示随机训练期望，不保证单次运行恰好达到。每步更新共享块及被选任务的监督路径，其他私有头不因此获得直接监督。目标系数相同仍可有不同抽取频率；1000步也不表示任务单步费用或容量改变后总计算相同。}
  \label{tab:ket-sharing-resource-controls}
\end{table*}

最后把资源和收益并到一起看，这也是和第15章评价的接口。三项任务的指标方向、量纲都不同：类别准确率、病害召回率越高越好，定位误差越低越好，绝不能把这三个原值直接加起来平均（高准确率和低误差相加没有意义）。正确顺序是：先逐项报告原始指标和各自的最低标准，再在明确定义的效用函数或对基线归一化之后，才汇总出任务平均、可兼容指标的样本平均、以及最弱任务的结果——多报告一个“最差任务”，是为了防止漂亮的均值掩盖某个任务的严重退化（呼应第11章均值方向反转的教训）。投入侧则要把标签数、独立对象数、重复抽取次数、训练算力、参数与优化器状态存储、以及部署时的调用条件全部列入；统一、严谨的评价要求由第~\ref{chap:ket-evaluation}章给出。

\begin{samepage}
\section{本章小结}
\label{sec:ket-sharing-summary}

全章可以收束为一句话：花卉的类别、病害与定位既有可以共享的共同线索，也有绝不能在共享中丢失的任务差异，而“能不能共享得好”从来不是某一个技巧决定的，它要求把四件事放在同一个方案里端到端检验——先依据交付责任选定哪些任务成组，再决定它们在哪一层交换表示或预测，然后确认损失加权与采样真的产生了你预期的那种更新，最后保证每项能力都分到了足够的训练样本和容量。还要破除一个幻觉：共享的参数或统一的接口只是“让交换成为可能”，它们本身并不保证有收益；共享到底有没有帮助，仍要回到相同评价条件、相同预算下的“单任务基线 vs 分组方案”对照来确定。

\end{samepage}

面对具体情形要有不同的判断纪律：只交付类别时，属性这类辅助任务的去留，只由“类别任务的净收益和额外成本”决定，不看它自己多准；三项都交付时，任何一项退化都要单独承担责任、不能被平均掩盖；多个地区副本需要个性化时，则要允许模型在共同中心、共享专家之外保留局部的差异化部分。还有两组极易混淆的“负”，这里最后再钉一次：训练中观测到的梯度负内积只是当前一步的局部冲突，不等于长期负迁移；参数协方差里的负相关是一种可能有益的统计互补，也不是“任务受损”的指标。当一个多任务系统效果不佳时，正确的排查姿势是顺着本章的链条逐项核查——关系测量准不准、判别特征有没有在共享中被丢掉、损失尺度和优先级对不对、标签可见性有没有偏、采样和预算够不够，而不是把所有原因都归咎于“某个权重没调好”。

最后是系统的联动观：这四类设计相互耦合，改动任何一项都可能动摇其他项的结论。重新分了组，就得重新测量任务亲和性（旧的两两关系不一定还成立）；扩展了专家，就得重新统计门控使用分布并复核容量账；更改了采样方案，就得重新计算真正被优化的有效目标（$p_kw_k$ 变了）。无论怎么改，做收益比较时都要守住三条不变量：任务单位一致、样本用途边界不被破坏、成本口径统一。还应再次明确本章的边界：它检验的是“已经在场的多个任务共同形成经验”的效果；若问一个没参与训练的新任务能否受益，那属于迁移，要整任务留出并规定适应权限（第13章）；而能力的长期保留、以及对总体证据强度的统一要求，则分别回到模型维护（第10、11章）和第15章的统一评价。
